\chapter{Bidirectional readers \texorpdfstring{\(K_D\)}{K(D)} and \texorpdfstring{\(K_\Omega\)}{K(Omega)}: profiles, coincidences, towers \texorpdfstring{\(90/120\)}{90/120}, and operators of the trajectory families}
\label{chap:lectores-bidireccionales-torres-operadores}

\section{Readers \(K_D\) and \(K_\Omega\): dodecaphasic profiles, path domain and spectral codomain}
For each oriented route, definitions are established.
\[
 K_D=\left(D_{27}+D_{36},D_1,\frac{D_4-D_3}{2}\right),
 \qquad K_\Omega=(\Omega,54,P_6).
\]
The exact census of the 324 routes contains fourteen \(K_D\) profiles, nine
\(K_\Omega\) profiles, forty pairs, and eighteen coincidences, all with
value \((80,54,6)\). These figures describe readers, not material species.

\section{Semigroup and fiber operators}
The global heights form
\[
 \langle90,120\rangle=\{90,120\}\cup\{30r:r\ge6\},
\]
with algebra \(\Bbbk[X_{90},Y_{120}]/(X_{90}^4-Y_{120}^3)\). The convergent
tower coefficients refine the character and define positive operators on
the thirteen multisections.

\section{Electronic construction, tower transducer and route clock}
The central electron provides the rank-one case; the rest of the
multisections requires the fiber operator, the coupling, and the spectral
reading. The transducer separately preserves temporal length, memory,
carry, and harmonic height.

\section{Dodecaphase ledger and TPK memory \texorpdfstring{\(1+5+6\)}{1+5+6}}

The operational selector, the calendar, and the dodecaphase record fix the
memory decomposition used by the mass readers. The identity
\(1+5+6\) belongs to the temporal state of the TPK and is not reconstructed from the
spectral outputs.

\begingroup
\let\chapter\section
\let\section\subsection
\let\subsection\subsubsection
\let\subsubsection\paragraph
\let\clearpage\relax
\let\cleardoublepage\relax
\section{Temporal dodecaphasic system of the TPK}
\label{sec:ciclo-dodecafasico-tpk-rev7}
\index{TPK!dodecaphasic system}\index{periodicity of order 108}

\HMTClaim{HMT-R-TPK-GAMMA108}

\subsection{Route selection operator and transport to the dimensional observable}

The ternary operative phase function is
\[
M_{\rm ph}(r)=
\begin{cases}
 +1,&r\in\{1,4,7\},\\
 -1,&r\in\{2,5,8\},\\
 0,&r\in\{3,6,9\}.
\end{cases}
\]
Its three values are the scalar inputs with which
\(\operatorname{Sel}_t\) activates, respectively, the additive cursor, the
multiplicative cursor or the neutral counter. Stroboscopic sampling is
\[
 \operatorname{Obs}_3(a)=(a_3,a_6,a_9,\ldots).
\]
The selection \(\operatorname{Sel}_t[\,M_{\rm ph}(g(t))\,]\) decides which
evaluation operates; cardinal transport moves the state; and
\(\operatorname{Obs}_3\) selects which instants are published.

Cardinal transport acts through four translations.
\[
 T_d(i,j)=(i,j)+\nu_d\pmod9,
\]
\[
 \nu_N=(-1,0),\qquad \nu_E=(0,1),\qquad
 \nu_S=(1,0),\qquad \nu_O=(0,-1).
\]
The additive and multiplicative cursors preserve separate positions and
orientations. The TRIT coordinate, the function \(M_{\rm ph}\), the operator
\(\operatorname{Sel}_t\), the four translations and
\(\operatorname{Obs}_3\) are therefore objects of different types.

\subsection{Dodecaphasic Calendar and Register}

The topological phase kernel does not terminate at the update of a position
on \(\mathbb Z_9^2\). Its minimal temporal state preserves the order of the
events, the orientation, the return, and the memory of the transitions. For
an orbit of \(108\) steps, the ordered partition is introduced
\[
 E_m=\{9(m-1)+1,\ldots,9m\},
 \qquad m=1,\ldots,12,
\]
and its ordered support is denoted by
\[
 \Gamma_{108}=E_1\sqcup\cdots\sqcup E_{12},
 \qquad 108=12\cdot9.
\]
Equivalently,
\[
 \Theta_{108}=\mathbb Z/108\mathbb Z,\qquad
 \beta(\theta)=
 \left[\left\lfloor\frac{\widetilde\theta}{9}\right\rfloor\right]_{12},
 \qquad
 r(\theta)=1+(\theta\bmod9).
\]
The coordinate \(\beta\) selects one of twelve sectors, and \(r\) locates the
phase within its nonadic window.

\begin{proposition}[Non-splittable central extension of the observable clock]
\label{prop:extension-central-108-no-escindida-rev8}
Writing \(C_n=\mathbb Z/n\mathbb Z\), the observable calendar forms the exact
sequence
\[
 0\longrightarrow C_{12}
 \xrightarrow{\ \iota\ }C_{108}
 \xrightarrow{\ p\ }C_9
 \longrightarrow0,
 \qquad
 \iota([b]_{12})=[9b]_{108},\qquad
 p([\theta]_{108})=[\theta]_9.
\]
For the set-theoretic section \(s([r]_9)=[r]_{108}\), with
\(0\leq r<9\), the gluing is given by the carry cocycle
\[
 c(r,r')=
 \left[\left\lfloor\frac{r+r'}{9}\right\rfloor\right]_{12},
 \qquad
 s(r)+s(r')=s([r+r']_9)+\iota\bigl(c(r,r')\bigr).
\]
The extension does not split. For the trivial action,
\[
 H^2(C_9,C_{12})
 \simeq C_{12}/9C_{12}
 \simeq C_3,
\]
and \(c\) represents the generating class \([1]\).
\end{proposition}

\begin{proof}
The image of \(\iota\) is the subgroup of multiples of nine in
\(C_{108}\), which coincides with \(\ker p\); hence exactness. The cocycle
identity is the Euclidean division of \(r+r'\) by nine. If the central extension
were split, the middle group would be isomorphic to
\(C_{12}\times C_9\). However,
\[
 \exp(C_{108})=108,\qquad
 \exp(C_{12}\times C_9)=\operatorname{mcm}(12,9)=36,
\]
which rules it out. The cohomological identification for a cyclic group with
trivial action gives \(H^2(C_9,C_{12})\simeq C_{12}/9C_{12}\), and the unit
carry projects to the generator of that quotient.
\end{proof}

The nonsplitting types the return: nine iterations close the coordinate in
\(C_9\) and translate one unit in \(C_{12}\); one hundred eight iterations close both
coordinates of this clock. The calendar \(C_{108}\) is the temporal
projection of the enriched state \(\mathcal X_{\rm TPK}^{\rm enr}\):
orientation, sheet, boundary, cylinder, genealogy, and deep memory
remain in the latter and can
distinguish histories with the same observable calendar.

If \(g_0(\theta)=\theta\bmod9\), then
\[
 g_0(\theta+9)=g_0(\theta),\qquad
 \beta(\theta+9)=\beta(\theta)+1\pmod{12}.
\]
The witness
\[
 (g_0,r,\beta,30^\circ\beta)(49)=(4,5,5,150^\circ),
\]
\[
 (g_0,r,\beta,30^\circ\beta)(58)=(4,5,6,180^\circ)
\]
shows exactly what returns and what retains memory: the nonadic phase and the internal position coincide, whereas the dodecaphase sector advances by one. The regional gauge \(\operatorname{diag}_c\) of the regions of \(\pi\) is not this coordinate \(\beta\). Nor is there a morphism \(30^\circ\beta\to h\in\langle90,120\rangle\) here; in particular, \(150\notin\langle90,120\rangle\). The pair \(49/58\) is a carry witness, not a regional or spectral selector.

For an enriched history \(x\) of TPK, let \(\tau_m\) be the ordered
subroute over \(E_m\), \(\sigma_m\) its orientation, \(\kappa_m\) the carry
that crosses the block boundary, and \(\Lambda_m\) the accumulated state of the
record. The object
\[
 \mathcal R_{12}
 =\bigl((E_m,\tau_m,\sigma_m,\kappa_m,\Lambda_m)\bigr)_{m=1}^{12}
\]
is the oriented dodecaphase temporal system of TPK. It is not a fourth
theory added to APP, TRIT, and TPK: it is the internal temporal form by which
TPK orders an orbit, preserves its returns, and publishes a reversible
record. The canonical projections are distinguished
\[
 \mathcal X_{\rm TPK}^{\rm enr}
 \longrightarrow \mathcal R_{12}
 \longrightarrow \Theta_{108}.
\]
The first preserves \((\tau_m,\sigma_m,\kappa_m,\Lambda_m)\); the second
publishes only the observable sequence and cannot by itself reconstruct the
signed state. The enriched record is denoted by \(\mathcal R_{12}\), and the
sector rotation by \(C_{12}\).

The observable closure of the calendar satisfies
\[
 F_{\rm obs}^{108}=I.
\]
Positions and orientations are already restored after \(54\) iterations,
whereas \(\mathcal R_{12}\) distinguishes the temporal position within the
twelve sectors. Here \(F_{\rm obs}\) is the observable permutation, whereas
\(\mathcal G_9\) is the pruning within the extension and boundary relation.
Monodromy is the joint property
\[
 g(K+9)=g(K),\qquad v_{K+9}\ne v_K,
\]
not the isolated name of the pruning. Chronological closure returns the observable
calendar; nonadic return preserves the phase and changes the prefix, cylinder,
and memory.

Two words of twelve components are also distinguished:
\[
 U_{12}^{\rm cat}=U_6\Vert U_6,
 \qquad
 U_{12}^{\rm sgn}=\mathcal H_{12}(K).
\]
The former concatenates two emissions of length six. The latter is a signed coordinate of the enriched state and is reversibly linked to \(K\) by the dodecaphase transformation \(\mathcal H_{12}\). Neither their domains nor their functions are interchangeable; in particular, the concatenated word does not by itself construct the signed coordinate.

\begin{definition}[Complete record of the twelve events]
\label{def:registro-dodecafasico-rev7}
For a route of cifras \(d_1,\ldots,d_{108}\), be
\[
 \mathcal D_{\rm evt}=\{2,4,6,8\},\qquad
 \Sigma_{12}=(+,+,+,-,-,-,+,+,+,-,-,-).
\]
Its event-oriented word is
\[
 e_m=\Sigma_{12}(m)\,
 \#\{t\in E_{m+1}:d_t\in\mathcal D_{\rm evt}\},
 \qquad 0\le m<12.
\]
The two semicrowns are paired through
\[
 p_i=e_i+e_{i+6},\qquad
 a_i=e_i-e_{i+6},\qquad 0\le i<6.
\]
The raw count, the five differences relative to the sixth pair, and
the six antisymmetries are defined by
\[
 s_C=\sum_{i=0}^{5}p_i,\qquad
 M_i=p_i-p_5\quad(0\le i<5),\qquad
 A_6=(a_0,\ldots,a_5).
\]
The full record is
\[
 \mathscr L_{12}(e)
 =\bigl(s_C,M_0,\ldots,M_4,a_0,\ldots,a_5\bigr).
\]
\end{definition}

\begin{theorem}[Exact coordinates of the dodecaphasic memory]
\label{thm:memoria-1-5-6}
The map \(e\mapsto\mathscr L_{12}(e)\) is injective on its image. Its
inverse is given by
\[
 p_5=\frac{s_C-\sum_{i=0}^{4}M_i}{6},\qquad
 p_i=p_5+M_i\quad(0\le i<5),
\]
\[
 e_i=\frac{p_i+a_i}{2},\qquad
 e_{i+6}=\frac{p_i-a_i}{2}.
\]
In particular, the distribution of coordinates is
\[
 \boxed{1+5+6=12}.
\]
\end{theorem}

\begin{proof}
By definition,
\[
 s_C=\sum_{i=0}^{4}(p_5+M_i)+p_5
 =6p_5+\sum_{i=0}^{4}M_i.
\]
This determines \(p_5\) and then all \(p_i\). Adding and subtracting the
equations \(p_i=e_i+e_{i+6}\) and \(a_i=e_i-e_{i+6}\) recovers the twelve
events. Divisibility by six and the six parity conditions
characterize the integral image; neither a mass nor a unit intervenes.
\end{proof}

Division by six is not applied during the count and is not a primitive
coefficient. First the raw oriented count \(s_C\) is formed; then the
record of twelve events is reconstructed; only then does
\(s_C/6\) appear in the character. Equivalently,
\[
 W_\pm=6n_A\pm s_C
\]
are integers and the lattice factorization has Smith normal form
\((1,12)\). This precedence will be mandatory in the mass law.

The integers \(90\) and \(120\) reappear later in six constructions
of distinct types. To prevent the name of the cardinal from supplanting the
object, the following convention is fixed from now on:

\begin{enumerate}
 \item \(\mathcal E_{\rm dod}=(D_3,D_4,Q):\mathbb Z^{12}\to\mathbb Z^{25}\),
       integral extractor that reconstructs \(K\);
 \item \(D_{\rm raw}^{90,120}(q)\), diagnostic difference of two series in
       a fixed chart;
 \item \(\iota_{6,8}:\mathcal F_6\hookrightarrow\mathcal F_8\), combinatorial
       inclusion of cardinalities \(90\) and \(120\);
 \item \(K_{6\to8}=12\Delta S_{90}-\Delta S_{120}\), angular functional
       defined after fixing channels, series, and normalization;
 \item \(\mathfrak S=\langle90,120\rangle\), infinite semigroup of spectral
       heights;
 \item \(\varepsilon_{\rm res}^{\circ}\), residue of an exact angular identity,
       not a finite extraction from any of the five preceding objects.
\end{enumerate}

The partial coincidence of indices establishes neither equality of operators,
domains, codomains, nor probative force. In particular,
\(K_{6\to8}\) uses two subqueues of the hierarchy but is not the complete
hierarchy; and neither of those subqueues forms a block independent of the mass
law.

\begin{remark}[Canonical Separation of Temporal Objects]
The notation distinguishes the oriented record \(\mathcal R_{12}\), the
observable calendar \(\Theta_{108}\), the sectorial rotation \(C_{12}\), and
the transport orbit \(\Gamma_{108}\), each with its domain, its
codomain, and its law of composition.
\end{remark}

\endgroup

\begingroup
\let\chapter\section
\let\section\subsection
\let\subsection\subsubsection
\let\subsubsection\paragraph
\let\clearpage\relax
% Vista pública derivada de 73_rev2_electron_lectores_torres.tex: cuerpo exacto y único retítulo académico del lector temporal de 108 estados. El fragmento del handoff permanece byte-exacto.
\section{Complete construction of the beta electron: central route, signature, character, and governing mass}
\Needspace{7\baselineskip}
\subsection{Central route of the electron and sequence of genealogical records}
The electron occupies the unique center \((5,5)\), with active orientation \(++\):

\[
 \Gamma_e=(5,5,++),\qquad
 \tau_{12}=+++---+++---.
\]
Their displacement readers are

\[
 (D_1,D_3,D_4,D_{27},D_{36})=(54,56,68,48,32).
\]
They must not be merged:

\begin{enumerate}
\item the raw signature \((24,0,12,0,-12)\);
\item the even CT--108 signature \((-12,-4,12,-6,12)\);
\item the vector of twelve events;
\item the affine origin \(v_e=0\) of the relative character.
\end{enumerate}
Nor should \(\Gamma_e\) be confused with the word \(w_e=201101\) of the reader of Euler's constant \(e\).

\Needspace{7\baselineskip}
\subsection{Basal electron selector and transition to the beta ruling value}
The internal selector, constructed without using the observed mass, closes

\[
 -22=-\left|((\mathbb Z/9\mathbb Z)\times\mathbb F_3)\setminus\iota(R_\pi)\right|,
 \qquad
 +15=|R_\pi\times\mathbb F_3|.
\]
The reading of vacancies also contributes

\[
 V_e=\begin{pmatrix}21&22\\6&7\end{pmatrix},\qquad
 \det V_e=15.
\]
The sign negativo of 22 records deficit/cut; the positive of 15, retencion/arean oriented. The matrix is local to the electron and does not  extiende as selector universal.

The basal stage is

\[
 e_0=\frac{\sqrt3}{4}
 \left[\exp\!\left(\frac{\varphi}{\pi^2}\right)
 -22\alpha_{\rm HMT}^3\right](1+15\Delta_4),
\]
\[
 e_0=0.510998922488\ldots\ \mathrm{MeV}.
\]
The entire genealogies are

\[
 22=\frac{396}{18}=27-5=24-2,
 \qquad
 15=\frac{270}{18}=\binom62=\binom64.
\]
\Needspace{7\baselineskip}
\subsection{Bidirectional readers and beta}
For every route,

\[
 K_D(\Gamma)=\left(D_{27}+D_{36},D_1,\frac{D_4-D_3}{2}\right),
\]
\[
 K_\Omega(\Gamma)=(\Omega_{90,120},54,P_6).
\]
Across the 324 routes there appear 14 profiles \(K_D\), 9 profiles \(K_\Omega\), 40 joint pairs, and 18 coincidences. The only common value is

\[
 (80,54,6).
\]
This 80 comes from the bidirectional reader of the electronic route, not from period 80 of the Euler word. Therefore

\[
 \kappa_e=80-54\alpha_{\rm HMT}+6\Delta_4,
\]
\[
 \boxed{
 \mathfrak e_e^\beta=e_0R_{\rm act}^{\kappa_e}
 =0.5109989506900008086082571648\ldots\ \mathrm{MeV}.}
\]
This is the governing electronic output. The basal output does not replace it. Equality with the antiparticle mass proceeds from route conjugation and from the parity of the mass operator, not from adding a second numerical row.

\Needspace{7\baselineskip}
\section{Physical realization by spectral coupling}
\Needspace{7\baselineskip}
\subsection{Obstruction to the equivariant punctual selection of a non-central multisection}
Let \(F\subset\mathcal R_{324}\) be a route fiber compatible with a structural class. Its linear space is

\[
 \mathcal H_F=\bigoplus_{\Gamma\in F}\mathbb C e_\Gamma.
\]
The central character and the readers define diagonal operators:

\[
 \widehat{\mathcal A}_F^{(0)}e_\Gamma
 =\chi_{\rm full}(\sigma_\Gamma,\eta_\Gamma)e_\Gamma,
\]
\[
 \widehat K_F^r e_\Gamma=\kappa_r(\Gamma)e_\Gamma,
 \qquad r\in\{D,\Omega\}.
\]
The two--way correction is

\[
 \widehat{\mathcal M}_F^{\beta,r}
 =R_{\rm act}^{\widehat K_F^r/2}
  \widehat{\mathcal M}_F^{(0)}
  R_{\rm act}^{\widehat K_F^r/2}.
\]
When \(\dim\mathcal H_F>1\), these operators generally have several eigenvalues. Choosing one because it approximates a known mass would reverse causality. Nor is it correct to require a physical name to identify a unique cell: the noncentral species resides in a representation of the fiber.

\Needspace{7\baselineskip}
\subsection{Spectral functor of the physical realization: positive coupling kernel, persistent subspace, and spectral measure}
A physical preparation \(P\) provides representation data, not a mass:

\[
 \mathfrak r(P)=
 (q,\,j,\,\epsilon_{2\pi},\,\epsilon_{4\pi},\,
  \text{statistics},\,\text{generation},\,\text{flavor},\,
  \text{color orbit},\,\text{interaction channel}).
\]
Let \(\mathfrak M\) be a coupling compatible with that representation. A positive kernel is defined

\[
 K_{P,\mathfrak M}:F\times F\longrightarrow\mathbb C
\]
which satisfies:

\begin{enumerate}
\item \(K=K^*\) y \(K\ge0\);
\item \(K(\Gamma,\Gamma')=0\) if the routes violate charge, parity, colour, statistics, or the
\(P\) channel;

\item \(K\) commutes with the symmetries that the coupling does not break;
\item \(K(C_M\Gamma,C_M\Gamma')=\overline{K(\Gamma,\Gamma')}\) under conjugation;
\item its entries are computed from the genealogical ledger, boundary, memory, carry, and holonomy, without
external magnitudes.

\end{enumerate}
The coupled evolution operator is

\[
 \mathcal U_{P,\mathfrak M}
 =K_{P,\mathfrak M}^{1/2}\,\mathcal U_{\rm TPK}\,
  K_{P,\mathfrak M}^{1/2}.
\]
The persistent subspace is obtained through the spectral projection onto the modes that survive return and pruning:

\[
 \Pi_{P,\mathfrak M}^{\rm surv}
 =\mathbf1_{\mathcal S_{\rm surv}}(\mathcal U_{P,\mathfrak M}).
\]
The correct physical realization is the spectral functor.

\[
 \boxed{
 \mathcal S_{\rm phys}^{\rm spec}(P,\mathfrak M)
 =\left(
 \mathcal H_{P,\mathfrak M}^{\rm surv},
 \widehat{\mathcal M}_{P,\mathfrak M},
 \mu_{P,\mathfrak M}
 \right),}
\]
where

\[
 \mathcal H_{P,\mathfrak M}^{\rm surv}
 =\Pi_{P,\mathfrak M}^{\rm surv}\mathcal H_F,
\]
\[
 \widehat{\mathcal M}_{P,\mathfrak M}
 =\Pi_{P,\mathfrak M}^{\rm surv}
  \widehat{\mathcal M}_F
  \Pi_{P,\mathfrak M}^{\rm surv},
\]
and \(\mu_{P,\mathfrak M}\) is its spectral measure in the prepared state.

When the coupling is specified by an irreducible representation \(\lambda\) of a finite group \(G_P\), the corresponding projector is obtained without choosing a distinguished route:

\[
 P_{P,\lambda}
 =\frac{\dim\lambda}{|G_P|}
  \sum_{g\in G_P}\overline{\chi_\lambda(g)}\,U(g).
\]
The bidirectional output is then the compressed pair.

\[
 \mathbb M(P,\lambda)=
 \left(
 P_{P,\lambda}\widehat{\mathcal M}^{\beta,D}P_{P,\lambda},
 P_{P,\lambda}\widehat{\mathcal M}^{\beta,\Omega}P_{P,\lambda}
 \right).
\]
If both operators commute with the irreducible representation, Schur's lemma forces their scalarity on the block. If they do not commute, the correct result is the spectrum of the minimal invariant block.

\Needspace{7\baselineskip}
\subsection{Naturality of the Spectral Functor under Intertwiners of the TPK Dynamics, Coupling, and Mass Operator}
If \(V:F\to F'\) intertwines the TPK dynamics, the coupling kernel, and the mass operators, then

\[
 V\Pi_{P,\mathfrak M}^{\rm surv}
 =\Pi_{P',\mathfrak M'}^{\rm surv}V,
 \qquad
 V\widehat{\mathcal M}_{P,\mathfrak M}
 =\widehat{\mathcal M}_{P',\mathfrak M'}V.
\]
By functional calculus, the spectral measure is transported by \(V\). The output does not depend on a numbering of the cells or on the choice of a basis within an orbit. This is the condition that must replace every non-equivariant nominal selector.

\Needspace{7\baselineskip}
\subsection{Conditions for the existence of a scalar mass under the dimensional reader}
A scalar mass is determined in any of the following cases:

\begin{enumerate}
\item the persistent subspace has rank one;
\item an irreducible representation fixes a unique spectral projector;
\item the measurement protocol defines a normalized positive state
\(\omega_{P,\mathfrak M}\) on the algebra of observables;

\item a resonance defines an isolated pole of the resolvent.
\end{enumerate}
In the third case,

\[
 m(P,\mathfrak M)
 =\omega_{P,\mathfrak M}
  (\widehat{\mathcal M}_{P,\mathfrak M}).
\]
In the absence of one of these mechanisms, the scientific output is the operator, its spectrum, and its multiplicities. That output is not incomplete: it is the object determined by the dynamics.

\Needspace{7\baselineskip}
\subsection{Finite construction of the mass kernel by means of representation, persistence, and coupling filters, followed by irreducible normalization}
The kernel is obtained in four finite stages:

\begin{enumerate}
\item \textbf{Representation filter.} Retain the routes whose sector, charge branch,
color, holonomy \(2\pi/4\pi\), parity, and statistics match \(P\).

\item \textbf{Persistence filter.} The partial order formed by class return,
orbit return, cell return, prefix separation, carry debt, and decimal stability is used. The result may be a Pareto orbit; uniqueness is not imposed.

\item \textbf{Coupling filter.} Boundary incidence and the
interaction channel are evaluated. Decoupled routes receive zero weight.

\item \textbf{Normalization.} On each irreducible component the trace of the
kernel is normalized to one.

\end{enumerate}
The oriented atlas provides all data of the first two stages. The third requires each physical realization to declare its channel; the fourth is canonical once the irreducible components have been fixed.

\Needspace{7\baselineskip}
\section{Genealogical Transducer of Towers}
\Needspace{7\baselineskip}
\subsection{Pair of associated words for the same genealogical route}
Each \(\Gamma\in\mathcal R_{324}\) determines simultaneously a word temporal

\[
 w_\Gamma\in\mathbb F_3^{108}
\]
and a dodecaphonic word signed

\[
 \tau_\Gamma\in\{-1,0,+1\}^{12}.
\]
The first preserves displacements along CT--108; the second preserves the accumulated orientation of the twelve events. The two words are linked by the same genealogical record, but they are not interchangeable. The refinement of towers reads both and also preserves the sheet, carry, winding, and return memory.

\Needspace{7\baselineskip}
\subsection{Correspondence between harmonic height and temporal length}
Write

\[
 \mathfrak S=30\langle3,4\rangle
 =\{90,120\}\cup\{30r:r\ge6\}.
\]
For \(h=30r\), the associated temporal length is

\[
 \ell_h=9r.
\]
In particular,

\[
 90\longleftrightarrow27,\qquad
 120\longleftrightarrow36,\qquad
 360\longleftrightarrow108.
\]
This correspondence does not identify the tower with time: it intertwines the height of the harmonic operator with the length of the route segment that generates it.

\Needspace{7\baselineskip}
\subsection{Temporal cycle shift reader of 108 states}
We define the cyclic discrepancy

\[
 D_\Gamma(n)=
 \sum_{t=0}^{107}
 \mathbf1_{\{w_{\Gamma,t}\ne
 w_{\Gamma,t+n\!\!\!\pmod{108}}\}}.
\]
The relative coefficient of the temporal branch is

\[
 \boxed{\eta_{30r}^{D}(\Gamma)
 =D_\Gamma(9r)-D_{\Gamma_e}(9r).}
\]
Subtraction does not use a mass as an origin: it uses the central route as the affine origin of the relative character. It holds.

\[
 \eta_{90}^{D}+\eta_{120}^{D}
 =[D_\Gamma(27)+D_\Gamma(36)]-80,
\]
so that the first layer exactly prolongs the first component of the reader \(K_D\).

\Needspace{7\baselineskip}
\subsection{Dodecaphase Accumulation Reader}
For \(r\ge1\), let

\[
 A_\Gamma(r)=
 \sum_{j=0}^{11}
 \left(
 \sum_{u=0}^{r-1}
 \tau_{\Gamma,j+u\!\!\!\pmod{12}}
 \right)^2.
\]
The relative coefficient of the oriented branch is

\[
 \boxed{\eta_{30r}^{\Omega}(\Gamma)
 =A_\Gamma(r)-A_{\Gamma_e}(r).}
\]
Then

\[
 4\eta_{90}^{\Omega}-3\eta_{120}^{\Omega}
 =\Omega(\tau_\Gamma)-80,
\]
so this layer extends the first component of \(K_\Omega\). The two families form the bidirectional lift

\[
 \mathcal F_{\rm tower}^{2w}:
 \mathcal R_{324}\longrightarrow
 \mathcal E_{\mathfrak S}\times\mathcal E_{\mathfrak S},
 \qquad
 \Gamma\longmapsto
 \bigl(\eta^D(\Gamma),\eta^\Omega(\Gamma)\bigr).
\]
\Needspace{7\baselineskip}
\subsection{Finite determination of the infinite tower}
The periodicity of the temporal word gives

\[
 D_\Gamma(9(r+12))=D_\Gamma(9r).
\]
If

\[
 T_\Gamma=\sum_{j=0}^{11}\tau_{\Gamma,j},
 \qquad r=12q+s,\quad 0\le s<12,
\]
then

\[
 A_\Gamma(12q+s)
 =A_\Gamma(s)+(12q^2+2qs)T_\Gamma^2.
\]
Therefore, the entire tower is determined by twelve values of \(D_\Gamma\), twelve values of \(A_\Gamma\), and the integer \(T_\Gamma^2\). The infinitude of heights does not require an infinitude of independent data.

The same formula solves the genealogical diamond of height 360. If

\[
 h(a,b)=90a+120b,
\]
then

\[
 (a+4,b)\sim(a,b+3)
\]
because \(3(a+4)+4b=3a+4(b+3)\). The genealogy \((a,b)\) is conserved until memory and carry are applied; only afterwards does it descend to the common height.

\Needspace{7\baselineskip}
\subsection{Absolute convergence of the refined character on the towers \(90/120\)}
We have

\[
 |\eta_h^D|\le108,\qquad
 |\eta_{30r}^{\Omega}|\le24r^2.
\]
Since \(0<q_+<q_-<1\),

\[
 |s_{30r}|\le2q_-^{30r}.
\]
It follows from this.

\[
 \sum_{h\in\mathfrak S}|\eta_h^D s_h|<\infty,
 \qquad
 \sum_{h\in\mathfrak S}|\eta_h^\Omega s_h|<\infty.
\]
The refined character is well defined without truncating the tower.

\Needspace{7\baselineskip}
\subsection{Memory, carry, and cocycle}
The bare coefficient is enriched by means of

\[
 \mathbb F_{\rm tower}(\Gamma)_h
 =\bigl(\eta_h^D,\eta_h^\Omega,\rho_h,
 \mathbf w_h,c_h,N_h^{\rm surv}\bigr),
\]
where the truncated signature satisfies

\[
 \sigma_{\Gamma|\ell_h}
 =r(\rho_h)+\mathcal D\mathbf w_h,
 \qquad
 \mathcal D=\operatorname{diag}(2,6,2,2,4),
\]
and the carry of two residues is

\[
 c(\rho_1,\rho_2)=
 \mathcal D^{-1}
 \bigl[r(\rho_1)+r(\rho_2)-r(\rho_1+\rho_2)\bigr].
\]
The locality defect of a composition,

\[
 \delta_h(\Gamma,\Lambda)=
 F_h(\Gamma\star\Lambda)-F_h(\Gamma)-F_h(\Lambda),
\]
obeys the cocycle identity

\[
 \delta_h(\Gamma,\Lambda)
 +\delta_h(\Gamma\star\Lambda,\Xi)
 =\delta_h(\Lambda,\Xi)
 +\delta_h(\Gamma,\Lambda\star\Xi).
\]
Thus it is preserved why two paths with the same visible height can carry distinct stories.

\Needspace{7\baselineskip}
\subsection{Selection of the read by coupling}
The two branches are not averaged for convenience. The physical coupling \(a\) must determine a natural transformation

\[
 \Pi_a:
 \mathcal E_{\mathfrak S}\times\mathcal E_{\mathfrak S}
 \longrightarrow\mathcal E_{\mathfrak S}
\]
that commutes with truncations and symmetries, respects the carry cocycle, preserves the electronic origin, and is compatible with \(X_{90}^{4}=Y_{120}^{3}\). If the coupling exchanges the two readings, the only continuous, multiplicative, symmetric, and normalized closure on positive scalars is the geometric mean; in exponents,

\[
 L_{\rm sym}=\frac{L_D+L_\Omega}{2}.
\]
In the absence of that symmetry, the output is the joint spectral measure of \((L_D,L_\Omega)\), not an arbitrary mean.

\Needspace{7\baselineskip}
\subsection{Orbit-primitive coinductive refinement}
The preceding finite lift admits a prolongation. Let \(\mathscr X_\Gamma\) be the graph of admissible extensions of the route, \(U_\Gamma\) its transfer operator, and \(R_\Gamma\) the orientation involution. The traces

\[
 a_n(\Gamma)=\operatorname{Tr}(R_\Gamma U_\Gamma^n)
\]
and the Möbius inversion

\[
 p_n(\Gamma)=\frac1n\sum_{d\mid n}\mu(d)a_{n/d}(\Gamma)
\]
separate primitive orbits. The difference

\[
 p_h(\Gamma)-p_h(\Gamma_e)
\]
is a further refinement of the height coefficient \(h\), not a replacement for the two finite readings.

\Needspace{7\baselineskip}
\subsection{Historical formulation via extension automaton}
The signature \(\mathbb Z^5\) does not exhaust the history of a route. For each \(\Gamma\in\mathcal R_{324}\), the following are preserved:

\begin{itemize}
\item the 108 transitions of the genealogical register;
\item the twenty blocks and their prefixes;
\item the palabras \(w_6\);
\item the nine comparisons \(K\leftrightarrow K+9\);
\item the leaf, the carry, and the winding;
\item the first returns of class, orbit, cell, and kinematic state;
\item the accumulated ambiguity and the stabilized decimal triads;
\item the profiles \(K_D\) and \(K_\Omega\).
\end{itemize}
These data allow the construction of a tower transducer without opening the residual of a mass.

\Needspace{7\baselineskip}
\subsection{Local extension automaton}
Let \(\mathscr X_\Gamma\) be the finite set of local states compatible with the route. A state preserves

\[
 x=(t,g,B,p,\varepsilon,c,w),
\]
where \(t\) is discrete time, \(g\) the nonadic phase, \(B\) the block, \(p\) the prefix, \(\varepsilon\) the sheet, \(c\) the carry and \(w\) the winding. An edge \(x\to y\) exists when \(y\) is an admissible extension of \(x\), respects the genealogical record and passes the pruning \(G_9\). The table of 36.864 carry compositions provides the local addition law; the return tables fix the memory.

Let \(U_\Gamma\) be the normalized transfer operator of this graph and \(R_\Gamma\) the orientation involution. It is required that

\[
 U_{C_M\Gamma}=J U_\Gamma J^{-1},\qquad
 R_{C_M\Gamma}=-J R_\Gamma J^{-1}.
\]
\Needspace{7\baselineskip}
\subsection{Primitive cell oriented counts and construction of the electronic signature}
For \(n\ge1\), the oriented trace

\[
 a_n(\Gamma)=\operatorname{Tr}(R_\Gamma U_\Gamma^n)
\]
counts signed closed returns. The Möbius inversion separates primitive orbits:

\[
 p_n(\Gamma)
 =\frac1n\sum_{d\mid n}\mu(d)a_{n/d}(\Gamma).
\]
The electron is taken as the affine origin of the relative corrections and is defined

\[
 \boxed{
 \eta_h(\Gamma)=p_h(\Gamma)-p_h(\Gamma_e),
 \qquad h\in\mathfrak S.}
\]
Thus the map is obtained.

\[
 \mathcal F_{\rm tower}:
 \Gamma^{\rm enr}\longmapsto
 \eta(\Gamma)=(\eta_h(\Gamma))_{h\in\mathfrak S}.
\]
There are no species-specific parameters. The same automaton, the same involution, and the same Möbius inversion act on the 324 routes.

\Needspace{7\baselineskip}
\subsection{Convergence of primitive refinement}
Let \(N_\Gamma=\dim\mathscr X_\Gamma\). If \(U_\Gamma\) is a contraction,

\[
 |a_n(\Gamma)|\le N_\Gamma.
\]
The number of divisors of \(n\) grows sublinearly, so

\[
 |p_n(\Gamma)|
 \le\frac{N_\Gamma\tau(n)}n.
\]
Since \(0<q_+<q_-<1\),

\[
 |s_h|\le q_-^h+q_+^h\le2q_-^h.
\]
By comparison,

\[
 \sum_{h\in\mathfrak S}|\eta_h(\Gamma)s_h|<\infty
\]
uniformly on any finite collection of fibers. The refined character is thus well defined.

\Needspace{7\baselineskip}
\subsection{Electronic route conservation identities and uniqueness of the central fixed point}
The construction must simultaneously satisfy:

\begin{enumerate}
\item invariance under renumbering of states;
\item covariance under leaf conjugation;
\item compatibility with direct sums of components;
\item equality upon recomputation from the genealogical register and from the certified graph;
\item conservation of the origin \(\eta(\Gamma_e)=0\);
\item absolute insensitivity to a permutation of the external mass columns;
\item determination of the tails by a bound established before comparison.
\end{enumerate}
The exact representation of a given residual proves completeness of the codomain, but does not replace these seven proofs.

\Needspace{7\baselineskip}
\section{Route clock and absolute normalization}
\Needspace{7\baselineskip}
\subsection{Integer returns are insufficient to distinguish species}
In the atlas of 324 routes, the first cell return is 27, the kinematic return within CT--108 is 54, and the extended kinematic return is 216. Being common, these integers fix the global temporal architecture, but by themselves they cannot explain mass differences between species.

\Needspace{7\baselineskip}
\subsection{Temporal holonomy of the electronic route and memory return}
The specific information resides in the holonomy of the extension operator. Let \(\lambda_j(\Gamma)\) be a nontrivial eigenvalue of \(U_\Gamma\) on the persistent subspace. We write

\[
 \lambda_j(\Gamma)=r_j(\Gamma)e^{i\theta_j(\Gamma)}.
\]
The associated internal frequency is

\[
 \omega_j(\Gamma)
 =\frac{\theta_j(\Gamma)+2\pi w_j(\Gamma)}{108t_0},
\]
where the integer \(w_j\) is fixed by the recorded winding, not by a subsequent choice of branch. For a stable mode \(r_j=1\); for a resonance \(r_j<1\), and its attenuation contributes a width.

The ratio

\[
 \rho_{\omega,j}(\Gamma)
 =\frac{\omega_j(\Gamma)}{\omega_0}
\]
is dimensionless and can modify mass ratios without altering the common decade \(-34\). The energy and mass of the mode are

\[
 E_j(\Gamma)=\hbar_{\rm ret}\omega_j(\Gamma)
 \chi_{\rm full}(\Gamma)R_{\rm act}^{\kappa_j(\Gamma)},
\]
\[
 m_j(\Gamma)=E_j(\Gamma)c_{\rm int}^{-2}.
\]
This equation brings together action, clock, character, and reader without converting any of them into another. Its numerical application requires determining \(U_\Gamma\), fixing the winding, and declaring the dimensional chart.

\Needspace{7\baselineskip}


\endgroup

\begingroup
\let\chapter\section
\let\section\subsection
\let\subsection\subsubsection
\let\subsubsection\paragraph
\let\clearpage\relax
% Vista científica exacta materializada desde una rama condicional.
% Fuente: source/demostraciones_primarias/14_05b_extension_superviviente_caracter.tex; rama: HMTIncludeRouteLedgerBlock.
% No modifica el contenido matemático de la rama de origen.
\chapter{Surviving extension, dodecaphasic record, and integral signature}
\label{chap:extension-ruta-caracter}

The trajectory that is observed does not constitute the primary object of this construction. The primary object is an extension compatible at all depths of HMT transport; a finite route is one of its readings. This distinction makes it possible to formulate unambiguously the internal chain leading from an enriched history to the action character. It also determines which information may be forgotten without altering the output and which must remain in the register.

The purpose of the section is to gather into a single argument four pieces that had
been studied separately: the surviving extension, the intrinsic
memory, the oriented record, and the integral signature. The terminal result of
this section is the formal character
\[
 \Gamma_{108}^{\rm enr}
 \xrightarrow{\ L_{\rm tick}\ }
 \ZZ^5
 \xrightarrow{\ \chi_{\rm form}\ }
 \operatorname{Fun}(\mathcal P,\RR_{>0}),
\]
prior both to specialization of the coefficients and to every physical unit. The first map extracts the signature from the register read; the second converts it into a positive function on the parameter space \(\mathcal P\). Numerical specialization is performed only after the global coefficients have been constructed in the corresponding section. The metrological realization belongs to a still later dimensional chart. The map \(L_{\rm tick}\) is the integral reader of iteration and return of the enriched chronological register.

\section{Operational composition of the prolongation and torsional cocycle in the generation of the route record}

The law of masses begins before the exponential formula. APP contributes a unique
support of routes and its two sheets, sum and product. TRIT orients each
elementary iteration.
If
\[
 r_t=1+((t-1)\bmod9),\qquad m_t:=M_{\rm ph}(r_t),
\]
the phase selector \(M_{\rm ph}\) activates
\[
 \operatorname{op}_t=
 \begin{cases}
  \Sigma,&m_t=+1,\\
  \Pi,&m_t=-1,\\
  \varnothing\quad(\text{\rm retention or neutral counter}),&m_t=0.
 \end{cases}
\]
The temporal reading may interpret the third case as threshold or present,
but it does not convert it into a third arithmetic sheet. APP is not what alternates
its sheets. TPK jointly transports
position, orientation, sheet, phase, block, carry, prefix,
cylinder, genealogy, and memory. The enriched chronological register records
each transition ---position,
sheet, reading, carry, phase, block, and event---. The local boundary signature
is subsequently derived through \(\operatorname{Sig}_q(h,c,\lambda)\), and the global
mass signature is extracted from the complete register; neither is reconstructed
retrospectively from a mass.

The finite funnel feeding this story preserves four distinct types:
\[
 9^2\!\cdot9^2\!\cdot4\!\cdot4=104\,976
 \longrightarrow 468\ \text{emissions }U_6
 \longrightarrow 243\ \text{words }w_6
 \hookrightarrow \mathbb F_3^6,
 \qquad |\mathbb F_3^6|=729.
\]
The first figure counts initial conditions of the two-cursor engine; 468
counts decimal emissions, and 243 visible ternary words. The 324 seeds
of the chronological projection \(\Pi_{108}^{\rm cal}\) of a single cursor
likewise do not replace the 104\,976
conditions of the complete state. Each arrow therefore changes both domain and
operation. The numerical equality between the ambient space of \(729\) words and the
\(729\) suffixes of a filtering stage is related only through the declared
ternary encoding; it does not identify the two objects.

The elevation preserves as well its causal order:
\[
 w_6\xrightarrow{\,L_0\,}w_{12}
 \xrightarrow{\,L_0\,}w_{18}
 \xrightarrow{\,L_1\,}w_{24}
 \xrightarrow{\,L_1\,}w_{30}
 \longrightarrow R_{36}\longrightarrow\mathcal G_9
 \longrightarrow
 \text{surviving section}.
\]
The boundary \(R_{36}\) opens the prolongation. At each filtration stage,
\(\mathcal G_9\) exhaustively ranges over
the 729 subcylinders and retains the unique one compatible with every horizon; the
deep specular episode \(3\to1\) is another pruning and must not be confused with
that ordinary partition \(729\to1\). The nonadic return preserves the phase
but changes the block, cylinder, and memory. Neither target figures nor masses intervene
in this selection. At the dimensional boundary, a particle is an enriched persistent
section of the extension; the route is its observable history.
Its physical name is read thereafter and does not decide which branch survives.

Before opening the dimensional chronological ledger, the same genealogy retains a
published window of twenty sextets per channel and admits the first
Archimedean reading of twelve decimal triads. These two witnesses belong to the
prolongation of the state: they are not the twelve events of the mass ledger, nor do they
replace them. The events are constructed only subsequently, by reading the route on the
oriented dodecaphase system.

The dodecaphasic genealogy subsequently converts the twelve oriented events into a
record and only then into the signature
\[
 \gamma_x\longmapsto\mathcal L(\gamma_x),\qquad
 \sigma_{\rm reg}:=L_{\rm tick}(\gamma_x)
 =(n_A,s_C,k_{\Dfour},b_{120},b_{\zeta}),
\]
\[
 \sigma_{\rm reg}\longmapsto
 \chi_{\rm form}(\sigma_{\rm reg};\,\cdot).
\]
The public coordinates \(b_{120}:=\nu_{120}^{\rm ev}\) and
\(b_{\zeta}:=\nu_{270}\) are outputs of the chronological ledger, not metrological
corrections introduced at the end. The subsequent coefficients of the
tower semigroup have a different provenance:
\[
 N_{120}=b_{120}+\eta_{120},\qquad
 b_{\zeta}\zeta_{270}\not\equiv\eta_{270}s_{270},
\]
with \(\zeta_{270}=135\Dfour^2\) and
\(s_{270}=q_-^{270}-q_+^{270}\). The tower refinement is added after
the central character without erasing the provenance of any coordinate.

Three lateral outputs also retain their type. The null sector leaves the positive branch through its null reader and null chart; it is not the limit \(m\to0\) of a positive character. A virtual trajectory is a finite extension with a terminal obstruction and preserves a chronological register and carry without thereby acquiring an imaginary mass. Finally, the Witt shadow accompanies every stage of filtration and preserves the exceptional genealogy but selects neither the route nor the mass coefficients.

\begin{figure}[p]
\centering
\begin{tikzpicture}[
  font=\small,
  node distance=3.5mm,
  box/.style={draw,rounded corners=1.5mm,align=center,text width=12.8cm,
    inner xsep=6pt,inner ysep=4pt},
  hmt/.style={box,draw=hmtblue,fill=hmtlightblue},
  md/.style={box,draw=hmtviolet,fill=hmtlightviolet},
  act/.style={box,draw=hmtgreen,fill=hmtlightgreen},
  result/.style={box,draw=hmtgold,fill=hmtlightgold},
  arr/.style={-{Stealth[length=2.2mm]},line width=.7pt}]

  \node[hmt] (app) {\textbf{APP}: nonadic arithmetic torus, paths, and additive
    and multiplicative sheets. The complete specification contains
    \(104\,976\) dual-cursor seeds.};
  \node[hmt,below=of app] (trit) {\textbf{TRIT}: local orientation of each
    transition; the lift preserves residue, quotient, and carry.};
  \node[hmt,below=of trit] (tpk) {\textbf{TPK}: phase transport with memory;
    \(104\,976\to468\,U_6\to243\,w_6\subset\mathbb F_3^6\).};
  \node[hmt,below=of tpk] (surv) {Coinductive filtration and monodromy
    \(\mathcal G_9\): surviving trajectory with conserved genealogy.};
  \node[md,below=of surv] (route) {The particle trajectory produces,
    without consulting the mass, a dodecaphasic record and the signature
    \(\sigma=(n_A,s_C,k_{\Delta_4},b_{120},b_\zeta)\in\mathbb Z^5\).};
  \node[md,below=of route] (chi) {The central character
    \(\chi_0(\sigma)\) and the harmonic refinement
    \(\eta\in\mathbb Z^{\langle90,120\rangle}\) construct
    \(\chi_{\rm full}(\sigma,\eta)\in\mathbb R_{>0}\).};
  \node[act,below=of chi] (action) {Independent dimensional branch:
    \(\Sigma_H=\varphi\alpha_{\rm HMT}^{16}\to\hbar_{\rm int}\),
    \(\omega_0=2\pi/(108t_0)\),
    \(m_0=\hbar_{\rm int}\omega_0c_{\rm int}^{-2}\).};
  \node[result,below=of action] (mass) {Typed confluence:
    \(m_X^{\rm int}=m_0\,\chi_{\rm full}(\sigma_X,\eta_X)\).
    The common scale fixes the absolute normalization and cancels in
    \(m_X/m_Y\).};
  \node[result,below=of mass] (compare) {The unit chart publishes the magnitude.
    CODATA and PDG are consulted only in the final comparison column.};

  \foreach \a/\b in {app/trit,trit/tpk,tpk/surv,surv/route,route/chi,
    chi/action,action/mass,mass/compare}{\draw[arr] (\a)--(\b);}
\end{tikzpicture}

\caption{Derivation of a mass magnitude from an HMT trajectory.
The character supplies a dimensionless ratio; the action scale and the clock
supply the dimensional section. The external reference intervenes only after
applying the chart of units.}
\label{fig:derivacion-masa-rev9}
\end{figure}

\section{Projection of a surviving extension onto its observable route}

Let \((\mathcal S_m,\rho_{m+1,m})_{m\geq0}\) be the system of surviving HMT states
at depth \(m\), with truncation maps
\(\rho_{m+1,m}:\mathcal S_{m+1}\to\mathcal S_m\). We define
\[
 \Omega_{\HMT}^{\rm surv}
 :=\varprojlim_m\mathcal S_m
 =\left\{(x_m)_m:\rho_{m+1,m}(x_{m+1})=x_m\right\}.
\]
An element \(x\in\Omega_{\HMT}^{\rm surv}\) simultaneously preserves the
compatibility of all its truncations. A reading chart of depth
108 produces
\[
 \operatorname{sh}_{108}(x)=\gamma_x\in\Gamma_{108}^{\rm enr}.
\]
The notation \(\operatorname{sh}\) recalls that \(\gamma_x\) is an
observable shadow: it records a sequential presentation of the state, but the state is not
defined by retrospectively choosing that route.

In addition to the visible word, the enriched presentation preserves phase,
sheet, orientation, boundary, antisheet memory, and the predecessors needed to
evaluate the cocycles. Two words with the same five final integers may
be distinct presentations because those integers are computed after
applying signs and correlations.

\begin{definition}[Intrinsic memory and read record]
An enriched dodecaphase extension is a compatible class of
truncations. We denote by \(\mathscr M(x)\) the intrinsic memory that those
truncations preserve. Once the chart
\(\gamma_x=\operatorname{sh}_{108}(x)\) is fixed, its read record is
\[
 \mathcal L(\gamma_x)
 =\bigl(n_A,e_0,\ldots,e_{11},T_\zeta,C_\zeta\bigr),
\]
where \(n_A\) is the count of the declared action support, the \(e_m\) are
the twelve oriented events, and \(T_\zeta,C_\zeta\) are the total torsional memory
and its restriction to the corona. Compatibility requires these data to
be preserved by every truncation containing the steps used in their
definition.
\end{definition}

The extension already contains the memory that the route presents; the route does not create it.
Relative to the declared chart, the composition
\[
x\longmapsto\gamma_x
\longmapsto\mathcal L(\gamma_x)
\longmapsto L_{\rm tick}(\gamma_x)
\]
is a map. Invariance under another presentation chart would additionally
require proving compatibility of the corresponding transition
maps. In this precise sense, the register is not added to an already
constructed particle: it is a reading of the extension's memory.

\section{Integral Decomposition of \(12\) Events into the Sectors \(1+5+6\)}

Let \(d_1,\ldots,d_{108}\in\{1,\ldots,9\}\) be the word of an enriched
presentation. In the parity convention used by the extractor, the support of
action is
\[
 \mathcal D_{\rm act}=\{1,3,5,7\},
 \qquad
 n_A=\#\{t:d_t\in\mathcal D_{\rm act}\}.
\]
The choice of this support is a declared part of the chart; it must not
be confused with the TRIT classification of the steps. For even events one
uses
\[
 \mathcal D_{\rm evt}=\{2,4,6,8\},
 \qquad
 \Sigma_{12}=(+,+,+,-,-,-,+,+,+,-,-,-).
\]
The oriented event of block \(m\in\{0,\ldots,11\}\) is
\[
 e_m=\Sigma_{12}(m)
 \#\{t\in\{9m+1,\ldots,9m+9\}:d_t\in\mathcal D_{\rm evt}\}.
\]

Let us pair the two semicrowns by means of
\[
 p_i=e_i+e_{i+6},
 \qquad
 a_i=e_i-e_{i+6},
 \qquad 0\leq i<6.
\]
We define
\[
 s_C=\sum_{i=0}^{5}p_i,
 \qquad
 M_i=p_i-p_5\quad(0\leq i<5),
\]
and write \(M_5=(M_0,\ldots,M_4)\),
\(A_6=(a_0,\ldots,a_5)\).

The \cref{thm:memoria-1-5-6}, in the primary residence of the
dodecaphasic TPK system, proves that
\[
 (e_0,\ldots,e_{11})\longmapsto(s_C,M_5,A_6)
\]
is injective over its image and provides the inverse
\[
 p_5=\frac{s_C-\sum_{i=0}^{4}M_i}{6},\qquad
 p_i=p_5+M_i,
\]
\[
 e_i=\frac{p_i+a_i}{2},\qquad
 e_{i+6}=\frac{p_i-a_i}{2}.
\]
Here only those coordinates are recovered to construct the character; the
proof and identity \(1+5+6=12\) are not duplicated.

The datum \(s_C\) is the one that enters the angular character. The eleven coordinates
\((M_5,A_6)\) preserve the placement of that total. Removing them before
composing histories conflates different events that may produce
different correlations.

\section{Torsional cocycle at the step level}

The event count does not by itself determine the torsional contribution. For
each occurrence of nine, retain the digit preceding it and define
\[
 j_2(d)=
 \begin{cases}
 0,&d=9,\\
 +1,&d\in\{2,4,6,8\},\\
 -1,&d\in\{1,3,5,7\}.
 \end{cases}
\]
With cyclic indices in the word, let it be
\[
 \zeta_t=\mathbf 1_{\{d_t=9\}}j_2(d_{t-1}),
\]
\[
 T_\zeta=\sum_{t=1}^{108}\zeta_t,
 \qquad
 C_\zeta=\sum_{t\in\{27,54,81,108\}}\zeta_t.
\]
The coordinate fuerte that acompana to \(\Dfour\) is
\[
 \boxed{k_{\Dfour}=T_\zeta-2C_\zeta.}
\]
The formula separates the accumulated drag from its oriented correction in the
four corona closures. The exhaustive census shows that neither
\(T_\zeta\) nor \(C_\zeta\) separately distinguishes all the census fibers;
the ordered pair does so in the two typed conventions being compared.

The two remaining tower coordinates are calculated from the event log. Let
\[
 \tau_m=\operatorname{sgn}(e_m).
\]
For the four step-four classes,
\[
 I_a=\{a,a+4,a+8\}\pmod{12},
 \qquad 0\leq a<4,
\]
we define the event-weighted character
\[
 \nu_{120}^{\rm ev}
 =\sum_{a=0}^{3}\operatorname{sgn}
 \left(\sum_{m\in I_a}e_m\right),
\]
and the three-step autocorrelation
\[
 \nu_{270}=\sum_{m=0}^{11}\tau_m\tau_{(m+3)\bmod12}.
\]
The variant
\[
 \nu_{120}^{\rm sgn}
 =\sum_{a=0}^{3}\operatorname{sgn}
 \left(\sum_{m\in I_a}\tau_m\right)
\]
is another observable. They are not silently identified: they differ on
25 of the 81 routes of the atlas.

\begin{definition}[Enhanced signature extractor]
Given the convention of events, it is defined
\[
 L_{\rm tick}(\gamma)=
 \bigl(n_A,s_C,k_{\Dfour},\nu_{120}^{\rm ev},\nu_{270}\bigr)
 \in\ZZ^5.
\]
The subscript indicates that the torsional coordinate is calculated before
the predecessors of the null steps are forgotten.
\end{definition}

This map is deterministic and consults neither masses, units, nor particle
labels. It is not asserted to be linear with respect to an already projected
concatenation: the signs and autocorrelations require the enriched
records to be composed first.

\endgroup
\begingroup
\let\chapter\section
\let\section\subsection
\let\subsection\subsubsection
\let\subsubsection\paragraph
\let\clearpage\relax
% Vista científica exacta materializada desde una rama condicional.
% Fuente: source/demostraciones_primarias/14_05b_extension_superviviente_caracter.tex; rama: HMTIncludeCharacterBlock.
% No modifica el contenido matemático de la rama de origen.
\chapter{Positive character of the signature and composition of records}
\section{Universal formal character}

Let \(\mathcal P=\RR^5\), with coordinates
\[
 X=(X_A,X_C,X_4,X_{120},X_{270}).
\]
For \(v=(n_A,s_C,k,\nu_{120},\nu_{270})\in\ZZ^5\), define
\[
 \ell_{\rm form}(v;X)
 =n_AX_A+s_CX_C+kX_4
  +\nu_{120}X_{120}+\nu_{270}X_{270},
\]
\[
 \chi_{\rm form}(v;X)=\exp\!\bigl(\ell_{\rm form}(v;X)\bigr).
\]
Thus, \(\chi_{\rm form}\) takes values in the multiplicative group
\(\operatorname{Fun}(\mathcal P,\RR_{>0})\), with pointwise product. It has not yet
selected global coefficients or a species.

\begin{theorem}[Universal Character of the Signature]
\label{thm:caracter-formal-ruta}
The map
\[
 \chi_{\rm form}:(\ZZ^5,+)\longrightarrow
 \bigl(\operatorname{Fun}(\mathcal P,\RR_{>0}),\cdot\bigr)
\]
is a homomorphism. The composition
\[
 \chi_{\rm route}^{\rm form}
 :=\chi_{\rm form}\circ L_{\rm tick}
\]
formally evaluates any enriched record and does not introduce a physical unit.
\end{theorem}

\begin{proof}
For every \(X\in\mathcal P\), the exponent
\(\ell_{\rm form}(\,\cdot\,;X)\) is additive on \(\ZZ^5\). Therefore,
\[
\chi_{\rm form}(v+w;X)
=\chi_{\rm form}(v;X)\chi_{\rm form}(w;X),
\]
and the identity holds as equality of positive functions.
\end{proof}

The formal character precedes the construction of the global coefficients.
This section does not yet define a numerical mass character or a
unit chart.

\section{Minimum internal section over the 81 cells}

The finite atlas contains 81 oriented cells. When they are projected onto the chart
of route family, 56 values are obtained. Denote this projection by
\[
 \mathcal R:\mathcal C_9\longrightarrow\mathcal R_{56}.
\]
The exact histogram of its fibers is
\[
 41\times1,
 \qquad10\times2,
 \qquad2\times3,
 \qquad1\times4,
 \qquad2\times5,
\]
and its sum is 81.

Having fixed the origin and the lexicographic order of the chart, let
\(\kappa(x)\in\{0,1,2,3,4\}\) be the rank of \(x\) within its fiber.

\begin{theorem}[Minimal internal selector]
\label{thm:selector-interno-minimo}
The map
\[
 x\longmapsto\bigl(\mathcal R(x),\kappa(x)\bigr)
\]
is injective on the 81 cells. No memory alphabet with fewer than
five symbols can lift \(\mathcal R\) injectively over the entire
atlas.
\end{theorem}

\begin{proof}
Within each fiber, \(\kappa\) assigns distinct ranks; between fibers, the
first coordinate is already distinct. There is a fiber of size five, so
the pigeonhole principle forces every injective lift to have at least
five memory values. The range \(0,\ldots,4\) attains that bound.
\end{proof}

Canonicity is relative to the declared origin and order. \(\kappa\) is neither a charge, a flavor, nor a mass. The thirteen coarse classes of the atlas produce finite multisections; only the central class possesses the fixed point \((5,5)\) under central inversion. The other twelve require oriented information in order to select a point and do not admit an equivariant point selector from the coarse class alone.

\subsection{Möbius Inversion and Round-Trip Characters}

On the word of dodecaphasic signs we define
\[
 C_M(\tau)_m=-\tau_{11-m}.
\]
This operation combines reversal of the reading order with the change of
sheet. A linear character
\(\ell_a(\tau)=\sum_ma_m\tau_m\), with
\(a_m=a_{11-m}\), satisfies
\[
 \ell_a(C_M\tau)=-\ell_a(\tau),
\]
while every symmetric quadratic form compatible with reversal
satisfies
\[
 Q(C_M\tau)=Q(\tau).
\]
Thus, the linear coordinates distinguish outward and return orientations, and the
quadratic coordinates preserve the correlation content. This parity explains
why \(s_C\) changes sign while \(\nu_{270}\) may remain
invariant. In the atlas, the seven word forms have multiplicities
\[
 54,2,2,2,7,7,7;
\]
therefore, the anti-sheet is realized as a correspondence between fibers, not as
a bijective involution obtained by ignoring their multiplicities.

\section{Orientation-based composition of records before projection}

Let \(\Gamma\) and \(\Lambda\) be two enriched histories, \(B\) a
common subhistory, and \(g\) the transport identifying phase, sheet,
orientation, and position at the interface. At ledger level, define
\[
 \boxed{
 \Gamma\star_B\Lambda
 =\bigl(n_\Gamma+n_\Lambda-n_B,
        e_\Gamma+g e_\Lambda-e_B\bigr).}
\]
The definition also presupposes identification of the edge, the
torsional predecessors, the corona, and the antisheet memory.

\begin{proposition}[Obstruction to the five-dimensional descent of the composition of records]
\label{prop:no-descenso-empalme-cinco}
In general, there is no binary operation on \(\ZZ^5\) that recovers
\(L_{\rm tick}(\Gamma\star_B\Lambda)\) from
\(L_{\rm tick}(\Gamma)\) and \(L_{\rm tick}(\Lambda)\) alone.
\end{proposition}

\begin{proof}
The operation composes the twelve events before applying
\(\operatorname{sgn}\) and the autocorrelations. Two histories may have
the same five-tuple and different values of \(M_5\) or \(A_6\). When composed
with a third history, some sum
\(e_a+e_{a+4}+e_{a+8}\) may cross zero in one but not the other, changing
\(\nu_{120}^{\rm ev}\). Finite diamonds supply explicit witnesses.
Therefore, projection onto five integers is not
a congruence for \(\star_B\).
\end{proof}

The proposition formalizes an essential operative rule: extensions are
composed; signatures are read afterward. The composite particle is likewise
not obtained by blindly adding two rows of a table, but by composing their
compatible memories and evaluating the result.

\section{Dodecaphasic dependence of the persistence character}

The reduced period 54 engines impose
\[
 e_{m+6}=e_m.
\]
It follows from here the congruences
\[
 n_A\equiv0\pmod2,
 \qquad
 s_C\equiv0\pmod2,
 \qquad
 \nu_{270}\equiv0\pmod4.
\]
The signatures of the abbreviated scalar reconstruction violate at least one of
these congruences. Therefore, their genealogy requires the enriched register and
does not arise from that reduced engine under the declared rules.

This result is not an obstruction to enriched TPK dynamics. It proves
what memory it requires: a genuine dodecaphase, with possible
\(A_6\ne0\), oriented composition, and retained torsional predecessors. The
ablation prevents a limitation belonging to a specific periodic reduction
from being attributed to TPK.

\section{Factorization by the quotient of log and physical realization}

In the scope of this section, the relation is declared.
\[
\gamma\sim_{\mathcal L}\gamma'
\quad\Longleftrightarrow\quad
\text{\(\gamma,\gamma'\) have the same extension and }
\mathcal L(\gamma)=\mathcal L(\gamma').
\]
Denote by
\(\mathcal E_{108}^{\rm surv}:=
\Gamma_{108}^{\rm enr}/\!\sim_{\mathcal L}\)
the quotient relative to this chart. Let
\[
 \pi_{\rm ext}:\Gamma_{108}^{\rm enr}\longrightarrow
 \mathcal E_{108}^{\rm surv}
\]
the canonical projection. By construction, the extractor is constant on its
fibers. The stronger assertion that every admissible chart of the same
extension automatically produces the same record would require proving
compatibility among charts and is not presupposed here.

\begin{theorem}[Relative factorization by the quotient of the logarithm]
\label{thm:descenso-caracter-extension}
There exists a unique application
\[
 \overline L:\mathcal E_{108}^{\rm surv}\longrightarrow\ZZ^5
\]
such that
\[
 \boxed{L_{\rm tick}=\overline L\circ\pi_{\rm ext}.}
\]
Therefore, \(\chi_{\rm form}\circ\overline L\) is a function of the class in
the relative quotient, not of the route representative chosen to read it.
\end{theorem}

\begin{proof}
If \(\pi_{\rm ext}(\gamma)=\pi_{\rm ext}(\gamma')\), both presentations
have the same record by the declared relation and, by the preceding
formulae, the same five-tuple. We define
\(\overline L([\gamma])=L_{\rm tick}(\gamma)\). Constancy proves that the
definition is independent of the representative; surjectivity of the map
to the quotient proves uniqueness.
\end{proof}

The internal construction thus reaches up to
\[
 \Omega_{\HMT}^{\rm surv}
 \longrightarrow\mathcal E_{108}^{\rm surv}
 \xrightarrow{\ \overline L\ }\ZZ^5
 \xrightarrow{\ \chi_{\rm form}\ }
 \operatorname{Fun}(\mathcal P,\RR_{>0}).
\]
The content of the theorem is positive and complete in this domain: a
surviving extension determines its ledger class, the class determines a
unique signature, and the signature determines the formal character. The global
specialization and the complete hierarchy of towers are constructed in the following
chapters; comparisons with physical names and units are presented
subsequently, without altering this chain.

\endgroup
\begingroup
\let\chapter\section
\let\section\subsection
\let\subsection\subsubsection
\let\subsubsection\paragraph
\let\clearpage\relax
\chapter{Split algebra and exact factorization of the character}
\label{chap:factorizacion-split}
\index{split algebra}\index{character!factorization split}
\index{null cone}\index{channels $q_\pm$}

The formal signature of an extension and the angular channels follow from branches already constructed. In this section, both converge in a determinant factorization that separately preserves the radial scale and the relative orientation. The result belongs to the algebraic stratum: it does not by itself identify the null ideals with physical photons or the positive interior with a material particle.

\section{Conjugate angular coordinates and bidirectional spectral lattice}
\label{sec:reticula-angular-split}

The preceding causal chain fixes
\[
 A_{\rm rad}=\frac{\pi A}{180},\qquad
 \Cstarrad=\frac{\pi\Cstar}{180},
\]
and, only after obtaining the direct coordinate \(A\) and the unique
conjugate coordinate \(C^*\),
\[
 q_+=\exp[-(A_{\rm rad}+\Cstarrad)],
 \qquad
 q_-=\exp[-(A_{\rm rad}-\Cstarrad)].
\]
For an angular signature \((n_A,s_C)\), we introduce the two winding
numbers
\[
 W_+=6n_A+s_C,
 \qquad
 W_-=6n_A-s_C.
\]

\begin{theorem}[Angular lattice of index twelve]
\label{thm:reticula-two-way}
The map
\[
 \mathcal W:\ZZ^2\to\ZZ^2,
 \qquad
 \binom{n_A}{s_C}\longmapsto
 \binom{W_+}{W_-}
 =\begin{pmatrix}6&1\\6&-1\end{pmatrix}
 \binom{n_A}{s_C}
\]
is injective, and its image has index twelve. Its invariant factors are
\((1,12)\), and
\[
 \operatorname{im}\mathcal W
 =\{(u,v)\in\ZZ^2:u+v\equiv0\pmod{12}\}.
\]
Moreover,
\[
 \boxed{
 \exp\!\left(n_AA_{\rm rad}+\frac{s_C}{6}\Cstarrad\right)
 =\left(q_+^{-W_+}q_-^{-W_-}\right)^{1/12}.}
\]
\end{theorem}

\begin{proof}
The matrix has determinant \(-12\). Since the greatest common divisor of its
entries is one, its Smith normal form is
\(\operatorname{diag}(1,12)\). If
\((u,v)=(6n+s,6n-s)\), then \(u+v=12n\). Conversely, if
\(u+v\) is a multiple of twelve, then
\[
 n=\frac{u+v}{12},
 \qquad
 s=\frac{u-v}{2}
\]
are integers and recover \((u,v)\). Finally,
\[
 \frac{W_+(A_{\rm rad}+\Cstarrad)
       +W_-(A_{\rm rad}-\Cstarrad)}{12}
 =n_AA_{\rm rad}+\frac{s_C}{6}\Cstarrad,
\]
which proves the identity upon exponentiation.
\end{proof}

The involution \(s_C\mapsto-s_C\) interchanges
\(W_+\leftrightarrow W_-\) and \(q_+\leftrightarrow q_-\). The twelfth
root is the integral index of the image of \(\mathcal W\), not a
rounding. The coordinate \(C^*\) acts as the antisymmetric component of the lattice;
it is neither a decimal correction to a mass nor an output of
\(\Gamma_{6\to8}\).

\section{Real split algebra}
\label{sec:algebra-real-partida}

\begin{definition}[Real split algebra]
Let
\[
 \mathcal A_{\rm sp}
 :=\RR[\mathsf R]/(\mathsf R^2-1).
\]
Conjugation and the norm are defined, respectively, by
\[
 \overline{a+b\mathsf R}=a-b\mathsf R,
 \qquad
 N_{\rm sp}(z)=z\overline z=a^2-b^2.
\]
\end{definition}

Conjugate idempotents
\[
 P_\pm=\frac{1\pm\mathsf R}{2}
\]
satisfy
\[
 P_\pm^2=P_\pm,\qquad
 P_+P_-=0,\qquad
 P_++P_-=1,\qquad
 \overline{P_+}=P_-.
\]

\begin{proposition}[Split spectral decomposition]
\label{prop:descomposicion-espectral-partida}
Every \(z\in\mathcal A_{\rm sp}\) admits a unique decomposition
\[
 z=r_+P_++r_-P_-,
 \qquad r_\pm\in\RR,
\]
and its norm is
\[
 \boxed{N_{\rm sp}(z)=r_+r_-.}
\]
Each principal ideal \(\mathcal A_{\rm sp}P_\pm=\RR P_\pm\) is a null
line; in particular,
\[
 N_{\rm sp}(\lambda P_\pm)=0.
\]
\end{proposition}

\begin{proof}
If \(z=a+b\mathsf R\), the choice \(r_+=a+b\), \(r_-=a-b\) provides the
decomposition, whose uniqueness follows from the independence of
\(P_+\) and \(P_-\). Conjugation exchanges the two coefficients:
\(\overline z=r_-P_++r_+P_-\). The orthogonality of the idempotents gives
\[
 z\overline z=(r_+r_-)(P_++P_-)=r_+r_-.
\]
The statement about each ideal is obtained by nullifying one of the two coefficients.
\end{proof}

\begin{definition}[Causal cone partitioned]
\label{def:cono-causal-partido}
The positive cone, its null boundary, and its interior are
\[
 \mathcal C_{\rm sp}^{+}
 =\{r_+P_++r_-P_-:r_+,r_-\geq0\},
\]
\[
 \partial_{\rm null}\mathcal C_{\rm sp}^{+}
 =\RR_{\geq0}P_+\cup\RR_{\geq0}P_-,
\qquad
 \operatorname{int}\mathcal C_{\rm sp}^{+}
 =\{r_+P_++r_-P_-:r_+r_->0\}.
\]
\end{definition}

The boundary has a single active direction and zero norm. The interior requires
two simultaneous directions and positive norm. This stratification is exact;
its physical realization requires maps that preserve the pertinent norm, transport, and
dynamics.

\section{Relative transporter and representational realization conditions}
\label{sec:transportador-relativo-split-md}

The historical block \(B_0=7I+2\mathsf R\) has spectrum \(\{9,5\}\), and the
angular block \(B_1=AI+\Cstar\mathsf R\) has spectrum
\(\{A+\Cstar,A-\Cstar\}\). The
\cref{thm:no-go-entrelazador-bloques-split} proves that, with the current internal
values, every putative linear intertwiner
\(\Phi B_0=B_1\Phi\) is zero. The active object is the relative
multiplicative transporter already constructed, not an equivalence between
representations with disjoint spectra.

\section{Partial factorization of the angular character}
\label{sec:factorizacion-split-angular}

Let
\[
 \sigma=\frac{s_C}{6},\qquad
 w_\pm=\frac{n_A\pm\sigma}{2}=\frac{W_\pm}{12},
 \qquad
 r_\pm=q_\pm^{-w_\pm},
\]
and let us define
\[
 Z_v=r_+P_++r_-P_-.
\]

\begin{theorem}[Partial factorization of the angular character]
\label{thm:split-character}
For every angular signature \(v=(n_A,s_C)\in\ZZ^2\),
\[
 \boxed{
 N_{\rm sp}(Z_v)
 =\det Z_v
 =\exp\!\left(
 n_AA_{\rm rad}+\frac{s_C}{6}\Cstarrad
 \right).}
\]
\end{theorem}

\begin{proof}
By the definition of the channels,
\[
 \log r_+
 =\frac{n_A+\sigma}{2}(A_{\rm rad}+\Cstarrad),
 \qquad
 \log r_-
 =\frac{n_A-\sigma}{2}(A_{\rm rad}-\Cstarrad).
\]
Upon summation,
\[
 \log(r_+r_-)
 =n_AA_{\rm rad}+\sigma\Cstarrad.
\]
The \cref{prop:descomposicion-espectral-partida} gives
\(N_{\rm sp}(Z_v)=r_+r_-\), and the diagonal representation in the basis
\((P_+,P_-)\) gives the same expression for the determinant.
\end{proof}

The component not seen by the determinant preserves the relative orientation:
\[
 \frac12\log\frac{r_+}{r_-}
 =\frac12\left(
 n_A\Cstarrad+\frac{s_C}{6}A_{\rm rad}
 \right),
\]
\[
 Z_v=
 \exp\!\left[
 \frac12\left(n_AA_{\rm rad}+\frac{s_C}{6}\Cstarrad\right)
 \right]
 \exp\!\left[
 \frac12\left(n_A\Cstarrad+\frac{s_C}{6}A_{\rm rad}\right)
 \mathsf R
 \right].
\]
The first factor fixes the radial scale; the second records the relative
orientation. They are two coordinated readings of one and the same element.

\section{Central incorporation of the towers}
\label{sec:factorizacion-split-torres}

Once \(\Dfour\), \(s_{120}\), and
\(\zCorr=\zeta_{270}^{\rm corr}\) have been constructed globally, let
\[
 t(v)=
 k_{\Dfour}\Dfour
 +\nu_{120}s_{120}
 +\nu_{270}\zCorr,
 \qquad
 \widetilde Z_v=e^{t(v)/2}Z_v.
\]

\begin{corollary}[Complete character as original norm]
\label{cor:split-character-torres}
Having fixed the preceding global blocks,
\[
 \boxed{
 N_{\rm sp}(\widetilde Z_v)
 =\exp\!\left(
 n_AA_{\rm rad}+\frac{s_C}{6}\Cstarrad
 +k_{\Dfour}\Dfour+\nu_{120}s_{120}
 +\nu_{270}\zCorr
 \right).}
\]
\end{corollary}

\begin{proof}
In the two-dimensional representation, multiplication by the central scalar
\(e^{t(v)/2}\) multiplies the determinant by \(e^{t(v)}\). The
\cref{thm:split-character} is applied.
\end{proof}

The corollary constructs the declared split realization of the active character;
it does not classify all possible realizations or transform the towers into
independent proofs. Its hypotheses are the already fixed channels and global
blocks, and its negative control is immediate: omitting the central factor
\(1/2\) doubles the tower exponent; simultaneously activating both
idempotents with positive coefficients departs from the null boundary.

\endgroup
\begingroup
\let\chapter\section
\let\section\subsection
\let\subsection\subsubsection
\let\subsubsection\paragraph
\let\clearpage\relax
\chapter{Values, Completeness and Blocks of the Spectral Hierarchy}
\label{chap:torres}

The bilayer operator, its moments, and the semigroup
\(\mathfrak S=\langle90,120\rangle\) are defined only once in the
\cref{sec:jerarquia-espectral-infinita-rev7}. This chapter publishes values,
proves reconstructive completeness, and separates the global blocks. The
same integers \(90,120\) also occur as cardinalities in the
hexad--octad incidence, but that coincidence does not transport the incidence
\(\iota_{6,8}\) to the angular chart or convert its configurations into
mass moments.

\section{First evaluations of the hierarchy}

The indices
\[
 90,120,180,210,240,270,360,420
\]
are a selection of initial heights of \(\mathfrak S\), since
\[
 180=2\cdot90,\quad210=90+120,\quad240=2\cdot120,
\]
\[
 270=3\cdot90,\quad360=3\cdot120,
 \quad420=2(90+120).
\]
The moments \(s_n=q_-^n-q_+^n\) and the generating interface
\(\mathcal K_{\mathrm{ang}}^{90,120}=(s_{90},s_{120})\) are used here with the
already fixed convention; they are not redefined. All values are
determined by \(A\), \(\Cstar\), and the reader's quadratic recurrence.

\begin{table}[htbp]
\centering
\caption{Oriented moments of the global hierarchy.}
\label{tab:torres}
\begin{tabular}{@{}rr@{}}
\toprule
\(n\)&\(s_n\)\\
\midrule
90  & \(2.9516943928982796\times10^{-4}\)\\
120 & \(1.9683511250294992\times10^{-5}\)\\
180 & \(8.7345871869508335\times10^{-8}\)\\
210 & \(5.8181313057291163\times10^{-9}\)\\
240 & \(3.8754674969305336\times10^{-10}\)\\
270 & \(2.5814553607509555\times10^{-11}\)\\
360 & \(7.6293257871406076\times10^{-15}\)\\
420 & \(3.3850663628348497\times10^{-17}\)\\
\bottomrule
\end{tabular}
\end{table}

\begin{proposition}[Global Character of Moments]
The quantities \(s_n\) of \cref{tab:torres} are invariants of the fixed
angular chart. A species contributes integer coordinates that weight those
invariants, but does not modify \(q_\pm\) or introduce a tower of its own.
\end{proposition}
\begin{proof}
Each \(s_n\) is a function of \(q_-\) and \(q_+\). The theorem of inversion of
channels proves that the joint map
\((A,\Cstar)\mapsto(q_+,q_-)\) is injective. No species label appears
in these definitions.
\end{proof}

\begin{theorem}[Reconstructive completeness of the hierarchy]
\label{thm:completitud-reconstructiva-torres}
Let \(n_1<n_2<\cdots\) be the increasing enumeration without repetitions of
\(\mathfrak S\). For
\[
 a_k:=s_{n_k}=q_-^{n_k}-q_+^{n_k}>0,
\]
every positive ratio \(\rho\in\RR_{>0}\) admits an absolutely
convergent expansion in the hierarchy:
\[
 \boxed{\rho=
 \exp\!\left(\sum_{k\geq1}\nu_k a_k\right)}
 \qquad\text{for some sequence }(\nu_k)_{k\geq1}\subset\ZZ.
\]
\end{theorem}

\begin{proof}
As \(0<q_+<q_-<1\), it holds
\[
 0<a_k<q_-^{n_k}.
\]
The \(n_k\) are multiplos of \(30\) no menores that \(90\); therefore,
\[
 \sum_{k\geq1}a_k
 \leq\sum_{m\geq3}q_-^{30m}
 =\frac{q_-^{90}}{1-q_-^{30}}<\infty,
 \qquad a_k\longrightarrow0.
\]
Fix \(r_0=\log\rho\) and a tie-breaking rule for the nearest integer.
Define recursively
\[
 \nu_k=\operatorname{nint}\!\left(\frac{r_{k-1}}{a_k}\right),
 \qquad
 r_k=r_{k-1}-\nu_k a_k.
\]
The property of the nearest integer gives
\[
 |r_k|\leq\frac{a_k}{2}\longrightarrow0.
\]
The telescoping identity
\[
 r_0=\sum_{j=1}^{N}\nu_j a_j+r_N
\]
implies \(r_0=\sum_{j\geq1}\nu_j a_j\). The convergence is absolute, since
\[
 \begin{aligned}
 \sum_{k\geq1}|\nu_k|a_k
 &=\sum_{k\geq1}|r_{k-1}-r_k|\\
 &\leq |r_0|+\frac{a_1}{2}
 +\frac12\sum_{k\geq2}(a_{k-1}+a_k)<\infty.
 \end{aligned}
\]
By exponentiating, the statement is obtained.
\end{proof}

\begin{remark}[Scope and negative control]
The theorem is a reconstructive surjectivity of the tower language onto \(\RR_{>0}\): it proves resolution of the codomain, not physical selection. The algorithm receives \(r_0=\log\rho\); if \(\rho\) is a target magnitude, its coefficients are therefore a target-dependent reconstruction. Prospective provenance requires another arrow,
\[
 \text{species}\dashrightarrow\text{extension}
 \longrightarrow\text{route}\longrightarrow\text{record}
 \longrightarrow\text{signature},
\]
and does not follow from this expansion. The conclusion would fail if the \(a_k\) did not
tend to zero, if the bound on the residue ceased to hold, or if the error series
were not summable.
\end{remark}

\section{Global blocks of mass action}

The mass action will use the triple
\[
 \Dfour,\qquad s_{120},\qquad \zCorr=135\Dfour^2.
\]
Their values are calculated only once. The coordinates
\(k,\nu_{120},\nu_{270}\) of a signature specify how many times
each block enters the exponent; they do not change its value.

It is useful to distinguish two objects that share the index \(270\):
\[
 \boxed{\zCorr=135\Dfour^2
 =1.66922699663643\times10^{-6}}
\]
and
\[
 \boxed{\sSpec=q_-^{270}-q_+^{270}
 =2.58145536075096\times10^{-11}.}
\]
The first is a quadratic invariant constructed from \(\Dfour\); the
second is a spectral moment of the semigroup. The coincidence of the subscript does not
authorize their identification.

\section{Higher-order combination}

The global combination
\[
 \delta\Dfour^{\rm tower}
 =\frac{s_{180}}{24}-\frac{s_{240}}{54}
 -24s_{360}+\frac76s_{420}
\]
evaluates a
\[
 \delta\Dfour^{\rm tower}
 =3.632051471908759123889070967\times10^{-9}.
\]
The comparison coordinate used in the selection of the
coefficients is defined as the scalar data
\[
 \delta\Dfour
 :=3.632051471908972784661641720\times10^{-9}.
\]
Consequently, the residue of the combination with respect to that coordinate is
\[
 \delta\Dfour^{\rm tower}-\delta\Dfour
 =-2.1366077257075\times10^{-22}.
\]
The small quantity is therefore a difference between two objects and not the value
of the combination. All the summands belong to the same hierarchy, and their
coefficients are global. Because the comparison coordinate intervened in the
selection of the coefficients, this equality is reconstructive. The
combination is not added again to the electronic formula that already uses
\(\Dfour\), since that substitution would duplicate the same correction.

\endgroup
\begingroup
\let\chapter\section
\let\section\subsection
\let\subsection\subsubsection
\let\subsubsection\paragraph
\let\clearpage\relax
\section{Bilayer Operator and Infinite Hierarchy}
\label{sec:jerarquia-espectral-infinita-rev7}
\index{towers!infinite hierarchy}\index{spectral semigroup}

Let
\[
 x=A_{\rm rad},\qquad y=C^*_{\rm rad},\qquad
 R_{\rm ch}^2=I,\qquad
 P_\pm^{\rm ch}=\frac{I\pm R_{\rm ch}}2.
\]
The reader of the two leaves meets in the operator
\[
 T_{\rm ang}=\exp(-xI-yR_{\rm ch})
 =q_+P_+^{\rm ch}+q_-P_-^{\rm ch},
\]
with
\[
 q_-=e^{-x+y},\qquad q_+=e^{-x-y}.
\]
Its even and odd spectral moments are
\[
 c_n=\operatorname{Tr}(T_{\rm ang}^n)=q_-^n+q_+^n,\qquad
 s_n=-\operatorname{Tr}(R_{\rm ch}T_{\rm ang}^n)=q_-^n-q_+^n.
\]
The involution \(C^*\mapsto-C^*\) preserves \(c_n\) and reverses \(s_n\).
Equivalently,
\[
 c_n=2e^{-nx}\cosh(ny),\qquad
 s_n=2e^{-nx}\sinh(ny).
\]
This is the convention fixed in the chapter on channels. To state without
homonymy the alternative presentation of the tower sources, set
\[
 R_{\rm tow}:=-R_{\rm ch},\qquad
 P_\pm^{\rm tow}=P_\mp^{\rm ch},
\]
\[
 T_{\rm tow}:=e^{-xI+yR_{\rm tow}}=T_{\rm ang},\qquad
 \operatorname{Tr}(R_{\rm tow}T_{\rm tow}^n)
 =-\operatorname{Tr}(R_{\rm ch}T_{\rm ang}^n)=s_n.
\]
The two writings describe the same scalar reader and do not constitute a
second orientation of the towers.

\begin{theorem}[Addition, norm, and recurrence]
\label{thm:recurrencia-bicapa-rev7}
For \(m,n\ge0\),
\[
 c_{m+n}=\frac{c_mc_n+s_ms_n}{2},\qquad
 s_{m+n}=\frac{s_mc_n+c_ms_n}{2},
\]
\[
 c_n^2-s_n^2=4e^{-2nx}.
\]
If \(p=q_-q_+=e^{-2x}\), both sequences satisfy
\[
 X_{n+2}=c_1X_{n+1}-pX_n.
\]
Moreover,
\[
 \partial_Ac_n=-\frac{n\pi}{180}c_n,
 \qquad
 \partial_As_n=-\frac{n\pi}{180}s_n,
\]
\[
 \partial_{C^*}c_n=\frac{n\pi}{180}s_n,
 \qquad
 \partial_{C^*}s_n=\frac{n\pi}{180}c_n.
\]
\end{theorem}

\begin{proof}
The addition identities follow by expanding
\((q_-^m\pm q_+^m)(q_-^n\pm q_+^n)\). The norm is
\[
 (q_-^n+q_+^n)^2-(q_-^n-q_+^n)^2
 =4(q_-q_+)^n.
\]
The last equality is the recurrence determined by the characteristic
roots \(q_-\) and \(q_+\). The differential equations are obtained
by differentiating the hyperbolic expressions; the factor \(\pi/180\) belongs to
the angular chart and not to a species selection.
\end{proof}

\begin{theorem}[Complete height semigroup]
\label{thm:semigrupo-alturas-rev7}
The heights of the transversal hierarchy form
\[
 \mathfrak S
 =\langle90,120\rangle
 =\{90a+120b:a,b\in\mathbb Z_{\ge0},\ a+b>0\}.
\]
Equivalently,
\[
 \boxed{\mathfrak S=\{90,120\}\cup\{30r:r\ge6\}.}
\]
\end{theorem}

\begin{proof}
Dividing by \(30\) yields the numerical semigroup
\(\langle3,4\rangle\). Its positive elements are
\(\{3,4\}\cup\{r:r\ge6\}\): \(5\) is the only positive gap beyond
the generators, and every \(r\ge6\) is obtained inductively by adding \(3\) or
\(4\). Multiplying again by \(30\) gives the formula.
\end{proof}

The presentation also preserves the genealogy of a height:
\[
 \mathfrak S\cong
 \bigl(\mathbb Z_{\ge0}^2\setminus\{(0,0)\}\bigr)
 \big/\bigl((a+4,b)\sim(a,b+3)\bigr),
\]
since \(4\cdot90=3\cdot120=360\). The height does not by itself determine the
composition that produced it; the genealogies \(4\cdot90\) and \(3\cdot120\)
remain distinct in the enriched chronological record even though they evaluate
to the same index.

To make explicit the recurrence sampled in the common step of thirty,
let us define
\[
 a_{30}:=q_-^{30},\qquad b_{30}:=q_+^{30},\qquad
 u_r:=s_{30r}=a_{30}^r-b_{30}^r.
\]
Then
\[
 \boxed{u_{r+2}=(a_{30}+b_{30})u_{r+1}-a_{30}b_{30}u_r.}
\]
The auxiliary term \(u_5=s_{150}\) may intervene when propagating the
recurrence, but \(150\notin\mathfrak S\); a computational identity does not create
a new admissible height.

In particular,
\[
90,120,180,210,240,270,300,330,360,390,420,450,480,\ldots
\]
are the first heights. The eight heights highlighted in previous
expositions are first readings of \(\mathfrak S\), not its definition nor
its limit.

As a human guide, its initial genealogy can be read without converting it into the
complete domain:
\[
\begin{array}{c@{\quad}c@{\quad}l}
90&3\cdot30&\text{minimal generator}\\
120&4\cdot30&\text{minimal generator}\\
180&2\cdot90&\text{first even mode}\\
210&90+120&\text{composite mode}\\
240&2\cdot120&\text{second even mode}\\
270&3\cdot90&\text{repetition of branch 90}\\
360&3\cdot120&\text{repetition of branch 120}\\
420&2(90+120)&\text{composite return}
\end{array}
\]
The last column records structural provenance; it neither introduces parameters
nor asserts that those eight rows exhaust the hierarchy. The same integers
\(90,120\) are incidence cardinalities in \(\mathcal F_6\) and
\(\mathcal F_8\), but that coincidence does not identify the bilayer operator with
the inclusion \(\iota_{6,8}\).

The interface of the two generators is only the pair
\[
 \mathcal K_{\rm ang}^{90,120}(q_+,q_-):=(s_{90},s_{120}).
\]
It must not be confused with the resolvents of two subtails or with the subsequent
functional \(K_{6\to8}=12\Delta S_{90}-\Delta S_{120}\), whose definition and
proof reside in \cref{mdg:chap:barbero}. There, the subtails live in the
analytic completion of the hierarchy; they are neither new generators nor
species coefficients.

\begin{definition}[Convergent module of spectral refinements]
\label{def:modulo-refinamientos-torre-rev7}
The integral domain of the global lift is
\[
 \mathcal E_{\mathfrak S}
 :=\left\{\eta\in\mathbb Z^{\mathfrak S}:
 \sum_{h\in\mathfrak S}|\eta_hs_h|<\infty\right\},
 \qquad s_h=q_-^h-q_+^h.
\]
Each \(\eta\in\mathcal E_{\mathfrak S}\) determines the spectral character
\[
 \mathcal R_\eta
 :=\exp\!\left(\sum_{h\in\mathfrak S}\eta_hs_h\right).
\]
\end{definition}

The mass law does not reside in this definition: its central signature
\((n_A,s_C,k_{\Delta_4},b_{120},b_\zeta)\in\mathbb Z^5\) is constructed beforehand
from the register. The elevation by \(\eta\) acts subsequently and preserves the
provenance of the basic term \(b_{120}s_{120}\). In particular,
\(N_{120}=b_{120}+\eta_{120}\) is the evaluated total coefficient, not a
new signature coordinate or a unique inverse decomposition.

The completeness of the semigroup does not by itself select a species. When
\(\eta\) is obtained from a supplied metrological residual, one has an
exact calibrated reconstruction; prospective force requires the route,
carry, and winding to produce \(\eta\) before that residual is opened. The
preceding recurrence generates all the moments once the
refinement has been fixed, without introducing a distinct law at each height.

\endgroup
\begingroup
\let\chapter\section
\let\section\subsection
\let\subsection\subsubsection
\let\subsubsection\paragraph
\let\clearpage\relax
\section{Pair of genealogical readers over the complete atlas}
\label{sec:operadores-masa-two-way-rev11}
\index{mass!bidirectional operators}
\index{atlas!chronological and dodecaphasic readers}

The enriched route possesses two complementary internal representations,
\[
 \Gamma\longmapsto
 \bigl(\gamma_{108}(\Gamma),\tau_{12}(\Gamma)\bigr)
 \in\mathbb F_3^{108}\times\{-1,0,+1\}^{12}.
\]
The first preserves the complete tritic history; the second preserves the
signed dodecaphase record. Each projection induces a natural reader, and both
are evaluated before opening masses, PDG labels, or residuals.

\HMTClaim{MD-R-MASS-TWOWAY-READER-PAIR}
\subsection{Pair of genealogical readers \(\mathbf K(\Gamma)=(K_D(\Gamma),K_\Omega(\Gamma))\): domain, codomain, and scalar evaluations}

Using the notation of
\cref{eq:KD-rev11,eq:kappaD-rev11,eq:kappa-omega-rev11}, we define
\begin{equation}
 \mathbf K(\Gamma)
 :=\bigl(K_D(\Gamma),K_\Omega(\Gamma)\bigr)
 \in\left(\mathbb Z^2\times\tfrac12\mathbb Z\right)
 \times\mathbb Z^3,
 \label{eq:lector-conjunto-masas-rev11}
\end{equation}
where
\[
 K_D=\left(D_{27}+D_{36},D_1,\frac{D_4-D_3}{2}\right),
 \qquad
 K_\Omega=(\Omega_{90\mid120},54,P_6).
\]
Scalar evaluations are
\begin{align}
 \kappa_D(\Gamma)
 &=D_{27}+D_{36}-\alpha_{\rm HMT}D_1
 +\frac{\Delta_4}{2}(D_4-D_3),
 \label{eq:kappaD-general-rev11}\\
\kappa_\Omega(\Gamma)
&=\Omega_{90\mid120}(\tau_\Gamma)
-54\alpha_{\rm HMT}+P_6(\tau_\Gamma)\Delta_4.
\label{eq:kappaOmega-general-rev11}
\end{align}
The second coordinate requires a causal distinction. In APP, the primordial spectral jump is \(9-5=4\); the local coefficient \(2\), the compression \(6=51-45\), and the two-dimensional unfolding \(54=9\cdot6=459-405\) are later typed invariants, in accordance with the \cref{thm:app-jerarquia-cuatro-dos-seis-cincuentaycuatro}. By contrast, in the chronological displacement reader, the electron coordinate is \(B_D=D_1(\Gamma_e)=54\): it counts immediate boundaries and coincides with the kinematic periodicity of order \(54\) of the CT108 endomorphism. The fixed \(54\) of \(K_\Omega\) is the common coordinate transferred when both readers are made compatible in the electron reduction. Its numerical equality with the global APP defect constitutes a subsequent cross-check; it does not make the number \(54\) the primordial difference between addition and multiplication.

The equality \(4\cdot90=3\cdot120=360\) explains the common genealogy of
the two readers. \(K_D\) compares displacements of the tritic word
\(\gamma_{108}\);
\(K_\Omega\) compares the energies of windows in the twelve-phase word. The
coincidence of origin does not identify their information quotients.
In the certified domain \(\Gamma_{324}^{\rm or}\), half-turn periodicity
makes all the \(D_k\) even, so that
\((D_4-D_3)/2\in\mathbb Z\) and \(K_D\in\mathbb Z^3\). Outside that domain,
the third component retains type \(\tfrac12\mathbb Z\).

\subsection{Naturalness of bidirectional operators and type preservation}

The reader \(D_k\) is invariant under time translation, reversal, and every
bijective permutation of the tritic alphabet. On the declared routes,
simultaneously reversing both orientations acts by a translation of
27 discrete steps and preserves \(K_D\). The reader \(K_\Omega\) is invariant under
dodecaphase rotation and reversal and under
\(C_M(\tau)=-\operatorname{rev}(\tau)\). Its quadratic form
\[
 \Omega(\tau)=\langle\tau,
 (4W_3^*W_3-3W_4^*W_4)\tau\rangle
\]
possesses integer bilinear polarization and satisfies the 2-cocycle identity.

The global counters of turn, sheet change, phase repetition, prefix change, and return are constant over the 324 routes. Consequently, every functional that factors exclusively through those six counters is constant and cancels in ratios. The readers in \eqref{eq:kappaD-general-rev11} and \eqref{eq:kappaOmega-general-rev11} act on the local history and pass that nonfactorization control.

\begin{ablation}[Necessity of the two sheets in the electronic route]
\label{abl:hojas-lector-two-way-rev11}
The chronological reader produces
\[
 (A_D,B_D,C_D)=
 \begin{cases}
 (80,54,6),&\substack{\text{joint chronological word of both sheets;}\\
 B_D=D_1(\Gamma_e)=54},\\
 (80,84,-8),&\text{additive projection},\\
 (80,60,13),&\text{multiplicative projection}.
 \end{cases}
\]
The complete electronic triple requires the joint history of the two arithmetic sheets.
\end{ablation}

\subsection{Exact Censuses over the 324 Routes}

The chronological reader produces fourteen profiles. The following table records their
complete census; \(A_D=D_{27}+D_{36}\), \(B_D=D_1\), and
\(C_D=(D_4-D_3)/2\).

\begin{longtable}{@{}rrr r@{}}
\caption{Exact census of the fourteen profiles of the chronological reader over the
324 oriented routes.}
\label{tab:censo-perfiles-kd-rev11}\\
\toprule
\(A_D\)&\(B_D\)&\(C_D\)&multiplicity\\
\midrule
\endfirsthead
\toprule
\(A_D\)&\(B_D\)&\(C_D\)&multiplicity\\
\midrule
\endhead
0&24&0&18\\
40&24&2&36\\
40&54&6&18\\
40&78&27&36\\
40&84&30&36\\
60&78&25&18\\
60&84&28&18\\
80&54&6&18\\
80&54&7&18\\
80&66&2&18\\
80&66&3&36\\
80&66&4&18\\
100&66&0&18\\
100&66&2&18\\
\midrule
\multicolumn{3}{r}{total}&324\\
\bottomrule
\end{longtable}

The dodecaphasic reader produces nine profiles:

\begin{table}[H]
\centering
\caption{Exact census of the nine profiles of the dodecaphasic reader over
the 324 oriented routes.}
\label{tab:censo-perfiles-komega-rev11}
\begin{tabular}{@{}rrr r@{}}
\toprule
\(\Omega\)&\(B_\Omega\)&\(P_6\)&multiplicity\\
\midrule
80&54&6&198\\
56&54&5&68\\
48&54&5&34\\
36&54&4&8\\
20&54&4&4\\
4&54&2&4\\
40&54&4&4\\
24&54&4&2\\
24&54&2&2\\
\midrule
\multicolumn{3}{r}{total}&324\\
\bottomrule
\end{tabular}
\end{table}

The readers coincide as triples on (18) routes, always with the value \((80,54,6)\), and the pair \(\mathbf K\) assumes (40) values. In the canonical section \(B1\) of 81 positions, \(K_D\) is constant in 41 of the 56 families and refines the quotient into 77 classes; \(K_\Omega\) is constant in 50 and produces 62 classes; the pair produces 78. In the lift to 324 orientations, \(K_D\) produces 146 classes and is nonconstant in every family; \(K_\Omega\) produces 100 and descends in 17 families; the pair produces 151 and is nonconstant in all of them. These censuses prove that the two projections retain distinct memories and determine the domain of each descent claim.

\subsection{Vector traces and joint spectra of the \(13\) multisections}

In the canonical section \(B1\), we denote by
\(\ell_j,\nu_j,d_j,u_j\) the generation-\(j\) multisections of the
charged-lepton, neutrino, down-type-quark, and up-type-quark
sectors, respectively. For each reader
\(r\in\{D,\Omega\}\) and each component \(a\in\{1,2,3\}\), we define the
diagonal operator
\[
 K_{M,a}^r e_\Gamma=K_{r,a}(\Gamma)e_\Gamma,
 \qquad
 \mathbf K_M^r=(K_{M,1}^r,K_{M,2}^r,K_{M,3}^r).
\]
The three operators of \(\mathbf K_M^r\) are self-adjoint and commute.
Therefore, their normalized trace is understood componentwise,
\[
 \overline{\mathbf K}_M^r
 =\frac{1}{\dim M}
 \bigl(\operatorname{Tr}K_{M,1}^r,
       \operatorname{Tr}K_{M,2}^r,
       \operatorname{Tr}K_{M,3}^r\bigr),
\]
and its spectrum is the joint spectrum, that is, the distribution with
multiplicity of the simultaneous triples. The thirteen vector traces are:

\begin{longtable}{@{}c r c c@{}}
\caption{Normalized vector traces of the readers \(K_D\) and
\(K_\Omega\) on the thirteen canonical multisections.}
\label{tab:trazas-multisecciones-two-way-rev11}\\
\toprule
\(M\)&\(\dim M\)&\(\overline{\mathbf K}_M^D\)&\(\overline{\mathbf K}_M^\Omega\)\\
\midrule
\endfirsthead
\toprule
\(M\)&\(\dim M\)&\(\overline{\mathbf K}_M^D\)&\(\overline{\mathbf K}_M^\Omega\)\\
\midrule
\endhead
\(\ell_0\)&1&\((80,54,6)\)&\((80,54,6)\)\\
\(\ell_2\)&8&\((65,249/4,17/2)\)&\((80,54,6)\)\\
\(\ell_4\)&16&\((225/4,489/8,21/2)\)&\((141/2,54,11/2)\)\\
\(\nu_1\)&2&\((40,51,29/2)\)&\((60,54,5)\)\\
\(\nu_2\)&4&\((45,63,29/2)\)&\((61,54,5)\)\\
\(\nu_3\)&6&\((170/3,55,34/3)\)&\((206/3,54,11/2)\)\\
\(\nu_4\)&8&\((105/2,237/4,95/8)\)&\((143/2,54,45/8)\)\\
\(d_1\)&2&\((40,51,29/2)\)&\((60,54,5)\)\\
\(d_2\)&4&\((45,63,57/4)\)&\((61,54,5)\)\\
\(d_3\)&6&\((170/3,55,23/2)\)&\((218/3,54,17/3)\)\\
\(d_4\)&8&\((105/2,237/4,49/4)\)&\((149/2,54,23/4)\)\\
\(u_1\)&4&\((65,66,9)\)&\((80,54,6)\)\\
\(u_3\)&12&\((205/3,119/2,113/12)\)&\((74,54,23/4)\)\\
\bottomrule
\end{longtable}

The notation \((a,b,c)^{\times n}\) indicates multiplicity \(n\). The complete
joint spectrum of \(\mathbf K_D\) is:

\begingroup\footnotesize
\begin{longtable}{@{}c p{.82\textwidth}@{}}
\caption{Complete joint spectra of the reader \(K_D\) in the thirteen
canonical multisections.}
\\
\toprule
\(M\)&\(\operatorname{Spec}_{\rm conj}\mathbf K_M^D\)\\
\midrule
\endhead
\(\ell_0\)&\((80,54,6)^{\times1}\)\\
\(\ell_2\)&\((0,24,0)^{\times1};(40,78,27)^{\times2};(80,54,7)^{\times1};(80,66,2)^{\times1};(80,66,3)^{\times1};(100,66,0)^{\times1};(100,66,2)^{\times1}\)\\
\(\ell_4\)&\((0,24,0)^{\times1}\); \((40,24,2)^{\times2}\); \((40,54,6)^{\times2}\); \((40,84,30)^{\times2}\); \((60,78,25)^{\times3}\); \((80,66,2)^{\times2}\); \((80,66,3)^{\times3}\); \((80,66,4)^{\times1}\)\\
\(\nu_1\)&\((40,24,2)^{\times1};(40,78,27)^{\times1}\)\\
\(\nu_2\)&\((0,24,0)^{\times1};(40,84,30)^{\times1};(60,78,25)^{\times1};(80,66,3)^{\times1}\)\\
\(\nu_3\)&\((40,24,2)^{\times2};(40,78,27)^{\times1};(60,84,28)^{\times1};(80,54,6)^{\times1};(80,66,3)^{\times1}\)\\
\(\nu_4\)&\((0,24,0)^{\times1};(40,24,2)^{\times1};(40,78,27)^{\times2};(40,84,30)^{\times1};(80,54,7)^{\times1};(80,66,2)^{\times1};(100,66,0)^{\times1}\)\\
\(d_1\)&\((40,24,2)^{\times1};(40,78,27)^{\times1}\)\\
\(d_2\)&\((0,24,0)^{\times1};(40,84,30)^{\times1};(60,78,25)^{\times1};(80,66,2)^{\times1}\)\\
\(d_3\)&\((40,24,2)^{\times2};(40,78,27)^{\times1};(60,84,28)^{\times1};(80,54,6)^{\times1};(80,66,4)^{\times1}\)\\
\(d_4\)&\((0,24,0)^{\times1};(40,24,2)^{\times1};(40,78,27)^{\times2};(40,84,30)^{\times1};(80,54,7)^{\times1};(80,66,3)^{\times1};(100,66,2)^{\times1}\)\\
\(u_1\)&\((40,54,6)^{\times1};(60,78,25)^{\times1};(80,66,2)^{\times1};(80,66,3)^{\times1}\)\\
\(u_3\)&\((40,24,2)^{\times2}\); \((40,78,27)^{\times2}\); \((60,84,28)^{\times1}\); \((80,54,6)^{\times3}\); \((80,66,3)^{\times1}\); \((80,66,4)^{\times1}\); \((100,66,0)^{\times1}\); \((100,66,2)^{\times1}\)\\
\bottomrule
\end{longtable}
\label{tab:espectros-kd-multisecciones-rev11}
\endgroup

The complete joint spectrum of \(\mathbf K_\Omega\) on the same fibres is:

\begingroup\footnotesize
\begin{longtable}{@{}c p{.82\textwidth}@{}}
\caption{Complete joint spectra of the reader \(K_\Omega\) in the thirteen
canonical multisections.}
\label{tab:espectros-komega-multisecciones-rev11}\\
\toprule
\(M\)&\(\operatorname{Spec}_{\rm conj}\mathbf K_M^\Omega\)\\
\midrule
\endfirsthead
\toprule
\(M\)&\(\operatorname{Spec}_{\rm conj}\mathbf K_M^\Omega\)\\
\midrule
\endhead
\(\ell_0\)&\((80,54,6)^{\times1}\)\\
\(\ell_2\)&\((80,54,6)^{\times8}\)\\
\(\ell_4\)&\((24,54,2)^{\times1};(56,54,5)^{\times4};(80,54,6)^{\times11}\)\\
\(\nu_1\)&\((40,54,4)^{\times1};(80,54,6)^{\times1}\)\\
\(\nu_2\)&\((4,54,2)^{\times1};(80,54,6)^{\times3}\)\\
\(\nu_3\)&\((36,54,4)^{\times1};(56,54,5)^{\times1};(80,54,6)^{\times4}\)\\
\(\nu_4\)&\((36,54,4)^{\times1};(56,54,5)^{\times1};(80,54,6)^{\times6}\)\\
\(d_1\)&\((40,54,4)^{\times1};(80,54,6)^{\times1}\)\\
\(d_2\)&\((4,54,2)^{\times1};(80,54,6)^{\times3}\)\\
\(d_3\)&\((36,54,4)^{\times1};(80,54,6)^{\times5}\)\\
\(d_4\)&\((36,54,4)^{\times1};(80,54,6)^{\times7}\)\\
\(u_1\)&\((80,54,6)^{\times4}\)\\
\(u_3\)&\((56,54,5)^{\times3};(80,54,6)^{\times9}\)\\
\bottomrule
\end{longtable}
\endgroup

The traces may coincide between multisections of different type; the type,
multiplicity, and full spectrum remain data of the operator.
The certificate reproduces these tables row by row.

\HMTClaim{MD-R-MASS-FIBRE-OPERATORS}
\subsection{Fiber Operators and Naturality of Base Change}

Distinguimos the fiber canonical \(F_M^{B1}\), formed by the positions of
the multisection \(M\) with its orientation published, and the fiber oriented
\(F_M^{\rm or}\), formed by their four lifts by position. For
\(s\in\{B1,{\rm or}\}\), let
\(\mathcal H_{M,s}=\bigoplus_{\Gamma\in F_M^s}\mathbb C e_\Gamma\). We define
\begin{equation}
 B_{M,s}^r e_\Gamma=\kappa_r(\Gamma)e_\Gamma,
 \qquad
 R_{\rm act}^{B_{M,s}^r}
 =\exp\bigl[(\log R_{\rm act})B_{M,s}^r\bigr],
 \qquad r\in\{D,\Omega\}.
 \label{eq:operadores-fibra-two-way-rev11}
\end{equation}
Each \(B_{M,s}^r\) is self-adjoint and diagonal in the route basis, so
\(R_{\rm act}^{B_{M,s}^r}\) is positive. For a positive basal mass
\(\widehat M_{M,s}^{(0)}\), the operatorial form conditioned on each reader is
\begin{equation}
 \widehat M_{M,s}^{\beta,r}
 =R_{\rm act}^{B_{M,s}^r/2}\,
  \widehat M_{M,s}^{(0)}\,
  R_{\rm act}^{B_{M,s}^r/2}.
 \label{eq:ley-operatoria-fibra-rev11}
\end{equation}
This expression is positive without imposing commutativity. When
\([\widehat M_{M,s}^{(0)},B_{M,s}^r]=0\), it reduces to
\(\widehat M_{M,s}^{(0)}R_{\rm act}^{B_{M,s}^r}\).
A unitary change of basis \(U\) transforms
\(B_{M,s}^r\mapsto UB_{M,s}^rU^*\) and preserves spectrum, trace, determinant, and
characteristic polynomial. The canonical electronic fiber \(F_e^{B1}\) has
rank one and reduces to the scalar of the
\cref{thm:reduccion-electron-two-way-rev11}.
Its oriented lift has rank four: the routes \((++)\) and \((--)\)
share the electronic profile, whereas the mixed orientations preserve
distinct profiles.

The normalized determinant
\[
 \operatorname{ndet}(R_{\rm act}^{B_{M,s}^r})
 =R_{\rm act}^{\operatorname{Tr}(B_{M,s}^r)/\dim F_M^s}
\]
is natural under permutation of the fiber. Logarithmic additivity distinguishes it as a geometric mean, but it is not used as a substitute for the physical realization. The exact refinement output remains the pair \((B_{M,s}^D,B_{M,s}^\Omega)\); the scalar section of a species is obtained through the flavor--generation--route morphism already constructed in \cref{eq:morfismo-sabor-ruta-cerrado-rev2,eq:realizacion-fisica-sabor-ruta-cerrada-rev2}: it selects the persistent route through non-metrological data and then applies the character of its signature. An arbitrary scalar functional \(S(B_{M,s}^D,B_{M,s}^\Omega)\) does not define that selection and remains only as a check that the target mass is not reintroduced.

\subsection{Absolute scale, relative correction and sensitivities}

The genealogy of the action line is written
\[
 \Sigma_H=\varphi\alpha_{\rm HMT}^{16},
 \qquad \operatorname{ord}_{10}(\Sigma_H)=-34,
\]
\[
 \delta_{2w}=H_5-\eta_{\rm ret}
 =\frac{90}{\pi}\mathscr C_\pi(\alpha_{\rm HMT}),
 \qquad R_{\rm act}=\frac{H_5}{\eta_{\rm ret}},
\]
\[
 \hbar_{\rm pre}=H_5\,10^{-34}\mathcal U_S,
 \qquad
 \hbar_{\rm ret}=\eta_{\rm ret}\,10^{-34}\mathcal U_S.
\]
The factor \(10^{-34}\mathcal U_S\) fixes the absolute section and cancels in
ratios. For the dodecaphase reader,
\begin{equation}
 \log\!\left[
 \frac{(m_X^{\beta,\Omega}/m_e^{\beta})}
 {(m_X^{(0)}/m_e^{(0)})}\right]
 =\bigl[(\Omega_X-80)+(P_{6,X}-6)\Delta_4\bigr]
 \log R_{\rm act}.
 \label{eq:razon-masa-omega-rev11}
\end{equation}
The common term \(-54\alpha_{\rm HMT}\) also cancels. For the chronological
reader,
\begin{equation}
 \log\!\left[
 \frac{(m_X^{\beta,D}/m_e^{\beta})}
 {(m_X^{(0)}/m_e^{(0)})}\right]
 =\bigl[\kappa_D(\Gamma_X)-\kappa_e\bigr]\log R_{\rm act}.
 \label{eq:razon-masa-D-rev11}
\end{equation}
With \(\log R_{\rm act}=6.932817628081925\ldots\times10^{-10}\),
\[
 \frac{\partial\log m}{\partial\Omega}=6.9328\times10^{-10},
 \quad
 \frac{\partial\log m}{\partial P_6}
 =\Delta_4\log R_{\rm act}=7.7090\times10^{-14},
\]
\[
 \frac{\partial\log m}{\partial\alpha}
 =-54\log R_{\rm act}=-3.74372\times10^{-8},
 \quad
 \left.\frac{\partial\log m}{\partial\Delta_4}\right|_e
 =6\log R_{\rm act}=4.15969\times10^{-9}.
\]
These derivatives are model sensitivities and remain separate from the
experimental uncertainty of a contrast mass.

\subsection{Physical pairing and null sector}

The typed pairing follows the chain
\[
 \begin{aligned}
 &\text{typed physical record}
 \longrightarrow\text{multisection or representation}
 \longrightarrow\text{fiber of routes}\longrightarrow\text{record}\\
 &\longrightarrow\text{signature and character}
 \longrightarrow(B_{M,s}^D,B_{M,s}^\Omega)
 \longrightarrow\text{dimensional chart}
 \longrightarrow\text{external comparison}.
 \end{aligned}
\]
The typed physical inventory distinguishes registers by object, representation, and codomain; the 324 routes are the internal engine; the 471 rows of the PDG catalogue are external properties. These cardinalities belong to different domains. The arrow is constructed from non-metrological physical descriptors, the multisection, and the route; a search by proximity to a mass belongs to the calibration regime.

The null sector bifurcates before the positive character. For photon and gluons,
\[
 \kappa_{\rm null}=\mathrm{N/A},
 \qquad m_{\rm null}=0.
\]
The expression \(R_{\rm act}^{0}=1\) is the unit of the positive character. Composite objects
concatenate records before evaluating the signature; topological sectors preserve fusion and
braiding; a virtual section is classified by its finite horizon of
prolongation.

\subsection{Independence of the path constructor with respect to masses and closure of the operator}

The structural verification uses exclusively the route census, the
register \(81\times108\), and the transducer of the 81 positions. It verifies their
row-by-row coupling, the coverage
\(81/56/13/324\), the naturality relations, the electronic triple \((80,54,6)\), the
censuses (14) and (9), the match (18/324), the
(40) pairs, differentiated descents in \(B1\) and in the 324 orientations,
periodicity 54, inversion by shift 27, and sheet ablations. The
null sector belongs to the domain exterior to the positive character. The calculation is
performed before introducing masses, CODATA, PDG, metrological residuals, or
nominal labels.

This certificate does not evaluate \(\alpha_{\rm HMT}\), \(\Delta_4\),
\(R_{\rm act}\), \(\kappa_e\), or the MeV chart. Those evaluations belong to
the arithmetic chain of the action line and to the energetic realization of
\cref{thm:reduccion-electron-two-way-rev11}; this arithmetic chain remains
separate from the combinatorial verification of the readers.

The metrological comparison subsequently applies the operators to the typed inventory,
declares unit, scheme, scale, date, and uncertainty, and performs permutations
and validation by retaining complete multisections. The physical output is
governed by the domain \(81/56/13/324\), by the two route readers, and by
the fiber operators.

\begin{statusbox}[title={Exact domain and closed physical realization}]
The atlas, the \(C_{81}\to\mathcal L_{108}\)-iteration chain of the
enriched chronological record of \(108\) iterations,
the action line,
the readers \(K_D,K_\Omega\), their electron reduction, and the fibre
operators are exact internal results. The pair
\(\mathbf K=(K_D,K_\Omega)\) is the complete output on each route and
determines the two diagonal operators of every fibre. In the electron
fibre, both readers coincide. In the noncentral fibres, the pair and its
spectrum preserve the complete refinement, while
\(\mathcal S_{\rm phys}\) binds branch, generation, fine signature, and, when
applicable, colour orbit to the persistent route. The character of that route
publishes the scalar coordinate before metrological comparison. Hence there is no
disjunction between a closed operator and a pending mass: they are two typed
readings of the same genealogy. Both readers are co-governing, and no
proximity to an external mass intervenes in the selection.
\end{statusbox}
\HMTClaim{MD-R-MASS-PHYSICAL-SELECTION-CLOSED}

\endgroup
\begingroup
\let\chapter\section
\let\section\subsection
\let\subsection\subsubsection
\let\subsubsection\paragraph
\let\clearpage\relax
% Vista editorial exacta de corpus/source/masas/04_persistencia.tex, líneas 365--590. Las líneas 1--364 ya residen exactamente en 73_rev2_electron_lectores_torres.tex.
\subsection{Persistence of the enriched state under the surviving extension and conservation of the route memory}
Let \(x\) be an APP state, \(\tau\) its tritic word, and \(B_K\) the memory block at level \(K\). The state relevant to mass has the form

\[
X_K=(x_K,\tau_K,\varepsilon_K,h_K,v_K,c_K,B_K,g_K),
\]
where \(\varepsilon_K\) is the sheet orientation, \((h_K,v_K)\) the cellular orientation, \(c_K\) the carry, and \(g_K\) the nonadic phase. An extension survives when it forms a compatible section of the successive cylinders. If \(C_{K,j}\), \(0\le j<729\) are the children of the cylinder at the next level, the pruning rule satisfies

\[
\sum_{j=0}^{728}G_9(C_{K,j})=1.
\]
The unique index \(j_K\) produces the transition

\[
P_{K+1}=729P_K+j_K.
\]
As nine levels are advanced, the phase is preserved,

\[
g(K+9)=g(K),
\]
but they change the block, prefix, and memory. Thus phase return is not state return. A stable particle is defined by the compatibility of the section across these returns, not by the motionless permanence of a visible coordinate.

The observable route \(\gamma_{\rm obs}\) is the chronological shadow of the extension. The point traversing a diagram is a reading cursor; the particle is the complete persistence class. This distinction permits a correct interpretation of quantum motion: what is required is not a small sphere following a classical trajectory, but a coherent history that admits several compatible projections.

\Needspace{7\baselineskip}
\subsection{Dodecaphasic registration and signature}
The temporal wheel contains twelve oriented events. Its causal decomposition is

\[
12=1+5+6:
\]
one event fixes the anchor, five determine the visible incidence, and six form the opposite pairs surviving the oriented quotient. If \(s_C^{\rm tot}\) is the integer contra-angular sum over the twelve events, the record preserves

\[
s_C:=s_C^{\rm tot}\in\mathbb Z,
\qquad
s_C^{(12)}:=\frac{s_C^{\rm tot}}6.
\]
The first object is the integer coordinate of the signature; the second is its dodecaphase angular reading. The divisor is not a normalization introduced from an external unit. It is the index that appears when the six opposite pairs are reconstructed. The dodecaphase incidence matrix has Smith normal form with factors \((1,12)\); the character is well defined on the quotient because the orientation residue is measured in that structure.

This is the governing convention of the entire work. The preserved historical expressions that designated the already divided coordinate as \(s_C\) are read as \(s_C^{(12)}\); their numerical values and attestations remain intact, but no current formula divides the same coordinate twice.

The direct winding pair is

\[
W_+=6n_A+s_C^{\rm tot},
\qquad
W_-=6n_A-s_C^{\rm tot}.
\]
As \(s_C=s_C^{\rm tot}\) in the published signature, this formula coincides with \(W_\pm=6n_A\pm s_C\). The contra-angular exponent uses \(s_C^{(12)}=s_C/6\) exactly once; it does not again divide a coordinate that had already been normalized.

The complete record publishes the signature

\[
\sigma(\Gamma)=
(n_A,s_C,k_{\Delta_4},b_{120},b_\zeta)\in\mathbb Z^5,
\qquad
\zeta_{270}=135\Delta_4^2.
\]
The five components have distinct types. \(n_A\) counts the direct angular development; \(s_C\) records the integer counter-angular sum; \(s_C^{(12)}=s_C/6\) is its angular reading; \(k_{\Delta_4}\) reads the global refinement; \(b_{120}\) marks the harmonic channel of height 120; \(b_\zeta\) belongs to the quadratic corrector of height 270. In particular, \(s_{270}\) and \(\zeta_{270}\) are not the same object: the former is an odd trace of the bilayer operator; the latter is a quadratic coordinate of the register.

The word electronic

\[
\tau_{12}=+++---+++---
\]
has linear sum zero, but its par signature is not zero:

\[
(Q_3,Q_4,Q_6,\rho_0,w)=(-12,-4,12,-6,12).
\]
Nor should it be identified with the raw signature of the electronic cycle
\((24,0,12,0,-12)\) or with the affine origin \(v_e=0\). These are four
different levels of reading: word, even signature, raw signature, and
origin of the relative character.

\Needspace{7\baselineskip}
\subsection{Angle, counter-angle, and character}
To make the construction autonomous, fix here the three internal coordinates that enter throughout the law:

\[
\alpha_{\rm HMT}
=\frac{7297352569283800997285105472380663}{10^{36}},
\]
\[
\Delta_4
=\frac{\pi_{\rm HMT}-e_{\rm HMT}\log\pi_{\rm HMT}}{270}
\left(1+\frac{\pi_{\rm HMT}}{729}\right),
\]
and, for \(a>0\) and \(\phi>0\),

\[
H_5(a,\phi)
=\frac12\sqrt{\frac{1000a}{\phi}}-a
+\frac{9a^2}{16}-\frac{5a^3}{9}
+\frac{7a^4}{48}-\frac{a^5}{54}.
\]
Henceforth,

\[
H_5:=H_5(\alpha_{\rm HMT},\varphi_{\rm HMT}).
\]
These definitions do not import an external metrological constant: \(\pi_{\rm HMT},e_{\rm HMT},\varphi_{\rm HMT}\) are the coinductive evaluations of the HMT reader, and \(\alpha_{\rm HMT}\) is the rational output of dodecaphase closure. The internal constant \(\alpha_{\rm HMT}\) precedes the counter-angle. The regional branch first determines

\[
A=1000\alpha_{\rm HMT},
\qquad
D_A=\exp\!\left(-\frac{100\pi\alpha_{\rm HMT}}9\right),
\]
\[
\mathcal C_\pi(\alpha_{\rm HMT})
=169\alpha_{\rm HMT}^{6}
+\frac{D_A\alpha_{\rm HMT}^{7}}
       {1-D_A\alpha_{\rm HMT}},
\qquad
y_*=\frac{\pi}{90}(H_5+\alpha_{\rm HMT})
-\mathcal C_\pi(\alpha_{\rm HMT}),
\]
\[
C^*=\frac{180}{\pi}y_*.
\]
Here \(A\) and \(C^*\) are already constructed angular coordinates, not action magnitudes. If

\[
\alpha_\angle:=\alpha_{\rm HMT}\,{}^\circ,
\qquad
\eta_{\rm ret}^{(\deg)}:=\frac{C^*}{2}-\alpha_\angle,
\]
then

\[
C^*=2\bigl(\eta_{\rm ret}^{(\deg)}+\alpha_\angle\bigr)
\]
is the inverse return identity, not a second rule for selecting the counter-angle. This notation avoids confusing the angular mantissa \(\eta_{\rm ret}^{(\deg)}\) with the dimensional magnitude \(\hbar_{\rm ret}\). The radian charts are obtained by multiplication by \(\pi/180\). The conjugate channels are

\[
q_\pm=\exp\left[-\frac\pi{180}(A\pm C^*)\right].
\]
The central character of a signature is

\[
\chi_0(\sigma)=
\exp\left(
n_AA_{\rm rad}
+\frac{s_C}{6}C^*_{\rm rad}
+k_{\Delta_4}\Delta_4
+b_\zeta\zeta_{270}
\right).
\]
The coordinate \(b_{120}\) remains in the signature, but its contribution is incorporated into the tower family. This convention prevents counting the same height twice.

The character is dimensionless. Its function is not to become an action constant, but to act by positive homotheties on an already typed mass line. For a pair of states \(X,Y\),

\[
\frac{m_Y^{(0)}}{m_X^{(0)}}
=\chi_{\rm full}(\sigma_Y-\sigma_X,\eta_Y-\eta_X).
\]
This formula separates the ratios basal of mass of the choice common of chart absolute. The quotient beta complete incorporates afterward \(R_{\rm act}^{\kappa_Y-\kappa_X}\).

\Needspace{7\baselineskip}
\subsection{Global mass tower semigroup and convergent harmonic expansion}
Let \(R^2=I\), with projectors

\[
P_\pm=\frac{I\pm R}{2},
\qquad
x=\frac{\pi A}{180},
\qquad
y=\frac{\pi C^*}{180},
\]
and

\[
T=e^{-xI-yR}=q_+P_++q_-P_-.
\]
The even and odd traces are

\[
c_n=\operatorname{Tr}(T^n)=q_-^n+q_+^n
=2e^{-nx}\cosh(ny),
\]
\[
s_n=-\operatorname{Tr}(RT^n)=q_-^n-q_+^n
=2e^{-nx}\sinh(ny).
\]
Satisfy

\[
c_n^2-s_n^2=4e^{-2nx},
\]
\[
c_{m+n}=\frac{c_mc_n+s_ms_n}{2},
\qquad
s_{m+n}=\frac{s_mc_n+c_ms_n}{2},
\]
and the second-order linear recurrence

\[
X_{n+2}=c_1X_{n+1}-e^{-2x}X_n,
\qquad X\in\{c,s\},
\qquad c_0=2,\quad s_0=0.
\]
The set of heights is not reduced to eight visible corrections. It is the semigroup

\[
\langle90,120\rangle
=\{90,120\}\cup\{30r:r\ge6\}.
\]
The heights \(90,120,180,210,240,270,360,420\) are the first useful genealogical witnesses, not the entire domain. The refined character is

\[
\log\chi_{\rm full}(\sigma,\eta)
=\log\chi_0(\sigma)
+\sum_{h\in\langle90,120\rangle}N_h(\sigma,\eta)s_h,
\]
with

\[
N_h(\sigma,\eta)=\eta_h+b_{120}\delta_{h,120}.
\]
Thus \(s_{120}\) appears only once, and at height 270, \(N_{270}s_{270}\) and \(b_\zeta\zeta_{270}\) remain separate. Absolute convergence requires

\[
\sum_h |N_hs_h|<\infty.
\]
The completeness of the semigroup guarantees the capacity to represent positive ratios; it does not, by itself, select a species' physical coefficients. The prospective selection rule, when used, must arise from the route, genealogical record, carry, memory, and winding before the external mass is consulted.

\Needspace{7\baselineskip}

\endgroup
\begingroup
\let\chapter\section
\let\section\subsection
\let\subsection\subsubsection
\let\subsubsection\paragraph
\let\clearpage\relax
% Vista editorial exacta de corpus/source/masas/06_operadores.tex, líneas 50--113. Las líneas 1--29 y 30--49 ya residen exactamente en 73_rev2_electron_lectores_torres.tex.
\subsection{Bidirectional readers \(K_D\) and \(K_\Omega\): profiles, coincidences, and operators by multisection}
The oriented route produces two families of invariants. For the temporal defects \(D_j\) and the return memory \(\Omega\), define

\[
K_D=\left(D_{27}+D_{36},D_1,\frac{D_4-D_3}{2}\right),
\qquad
K_\Omega=(\Omega,54,P_6).
\]
Across the 324 routes there are 14 distinct profiles of \(K_D\), 9 profiles of \(K_\Omega\), 40 joint pairs, and 18 coincidences. The unique profile shared by the electron route is

\[
K_D(\Gamma_e)=K_\Omega(\Gamma_e)=(80,54,6).
\]
These readers are not decorative labels. \(K_D\) measures discrepancies in the word; \(K_\Omega\) measures the enriched return. Their equality at the center certifies rank-one bidirectional closure. Away from the center, they produce fiber operators whose complete spectra must be published even when a subsequent physical realization selects a scalar section.

\Needspace{7\baselineskip}
\subsection{Proper harmonic refinement of each state}
The full refinement lives in the coefficient space.

\[
\mathcal E_{90,120}
=\left\{\eta=(\eta_h)_{h\in\langle90,120\rangle}:
\sum_h|\eta_hs_h|<\infty\right\}.
\]
The sought application does not receive a mass. It receives the enriched route:

\[
\mathfrak T:\Gamma_{\rm enr}\longrightarrow\mathcal E_{90,120},
\qquad
\Gamma\longmapsto\eta(\Gamma).
\]
Their coordinates must be explicit functions of returns, carry, winding, memory, and boundary. The naturality condition is

\[
\eta(r_{n,m}\Gamma)=r^{\mathcal E}_{n,m}\eta(\Gamma),
\]
and the convergence must be accompanied by a computable tail bound. The complete character then acts by

\[
\widehat M_F
=\sum_{\gamma\in F}
m_{\rm clk}\,
\chi_{\rm full}(\sigma(\gamma),\eta(\gamma))
|\gamma\rangle\langle\gamma|.
\]
The capacity to reconstruct a given numerical difference proves representational completeness of the hierarchy; it does not by itself prove that \(\mathfrak T\) has chosen its coefficients. The difference between the two assertions is preserved in every quantitative row.

\Needspace{7\baselineskip}
\subsection{Coupling Scalarization}
Let \(P_{X,a}\) be the projector determined by the state \(X\) and the apparatus or physical channel \(a\). The mass distribution is

\[
\mu_{X,a}(B)
=\langle\Psi_X,
P_{X,a}E_{\widehat M_F}(B)P_{X,a}\Psi_X\rangle.
\]
There is a scalar line when the compression satisfies

\[
P_{X,a}\widehat M_FP_{X,a}=m_XP_{X,a}.
\]
This is the precise meaning of “to measure is to couple.” The value was not hidden in a nominal cell; it appears as an eigenvalue of a physical compression of an already constructed operator. If the compression retains several eigenvalues, the correct prediction is a multiplet or a distribution, not an average chosen to match a reference.

\Needspace{7\baselineskip}


\endgroup
\begingroup
\let\chapter\section
\let\section\subsection
\let\subsection\subsubsection
\let\subsubsection\paragraph
\let\clearpage\relax
\subsection{Physical selection as a natural relation}
Away from the electronic center, a species must not select a single cell by decree. The twelve noncentral multisections are fibers stable under cellular inversion and do not, in general, admit an equivariant point section. The correct realization begins with a relation

\[
\mathscr S\subset \mathcal P\times\mathcal R_{324},
\qquad
\mathscr S(X)=\{\gamma:(X,\gamma)\in\mathscr S\},
\]
whose value may be a point, an orbit, or a fiber. For the electron, \(\mathscr S(e)=\{\gamma_e\}\) has rank one because \((5,5)\) is the unique fixed point. For color, flavor, and generation, \(\mathscr S(X)\) preserves the corresponding orbit until a physical coupling compresses it.

The relation is admissible if it satisfies simultaneously:

\begin{enumerate}
\item depends only on representation, orientation, memory, boundary, charge,
flavor, color, and generation constructed before the contrast;

\item commutes with particle--antiparticle conjugation;
\item is equivariant with respect to the chromatic action and cellular inversion;
\item is compatible with refinement, nonadic return, and dodecaphase reading;
\item preserves every multiplicity that has not been eliminated by a coupling;
\item remains invariant under permuting or replacing the external mass values.
\end{enumerate}
The sixth point is the decisive control. If altering the experimental column changes \(\mathscr S(X)\), the construction is retrospective. If the relation remains and only the contrast residual changes, the causal direction is preserved.

\Needspace{7\baselineskip}

\endgroup

\section{Scalarization morphism and dimensional realization obstructions}
Outside the rank-one electronic fibre, the spectrum does not itself select a
physical eigenvalue. A scalar requires a projector, state, or coupling functional
declared before external comparison.
